% !TEX root = ../arxiv.tex
\makeatletter
\let\addcontentsline\arxiv@savedaddcontentsline
\makeatother

\newpage
\setcounter{tocdepth}{2}
\renewcommand{\contentsname}{Appendix Contents}
\tableofcontents
\newpage
\appendix

\section{Supporting Definitions and Proofs for Section~\ref{sec:method:schedule:signal}}
\label{app:proofs-schedule}

This appendix provides the justification of timestep alignment in Spectral Progressive Diffusion, the deferred proofs of Propositions~1 and~2 from Sec.~\ref{sec:method:schedule:signal}, along with supporting definitions.

\subsection{Definitions}
\label{app:prelim-schedule}

\paragraph{Per-frequency SNR.}
Following~\citet{kingma2021vdm}, we define the \emph{per-frequency signal-to-noise ratio (SNR)} as the ratio of the clean- and noise-component second moments of the per-frequency forward process in Eq.~\eqref{eq:forward-mode}:
\begin{equation}
  \mathrm{SNR}_\omega(t)
  \;:=\;
  \frac{\mathbb{E}[\,|(1-t)\,x_0^{(\omega)}|^2\,]}{\mathbb{E}[\,|t\,\epsilon^{(\omega)}|^2\,]}
  \;=\;
  \underbrace{\frac{(1-t)^2}{t^2}}_{\mathrm{SNR}(t)}\;\cdot\;P_\omega,
  \label{eq:snr}
\end{equation}
which factorizes into a frequency-independent timestep term $\mathrm{SNR}(t) = (1-t)^2/t^2$ and the frequency-specific signal power $P_\omega$.

\paragraph{Resolution grids and representable frequencies.}
By the Nyquist--Shannon sampling theorem~\citep{nyquist1928certain,shannon1949communication}, a spatial grid of size $H \times W$ can represent 2D frequencies in the rectangle $|\omega_x| \le W/2$, $|\omega_y| \le H/2$. We define the maximum (fully) representable radial frequency as
\begin{equation}
  \omega_{\max}(H, W) \;=\; \frac{\min(H, W)}{2}.
  \label{eq:nyquist-cap}
\end{equation}
Note that our definition of the maximum (fully) representable frequency above is conservative, as it captures only the maximum radial frequency fully supported on the spatial grid $H \times W$.
A stage operating at reduced resolution $(sH, sW)$ for $s \in (0,1]$ can represent only a subset of the representable frequencies at full-resolution:
\begin{equation}
  \Omega_s \;:=\; \{\,\omega \in \Omega : |\omega| \le s \cdot \omega_{\max}(H, W)\,\} \;\subseteq\; \Omega.
  \label{eq:nyquist-set}
\end{equation}

\subsection{Justification of Timestep Alignment}
\label{app:time-align}

We verify that applying the scalar rescaling in Eq.~\eqref{eq:scale-align} and querying the pretrained model at the aligned timestep $\tilde{t}_i$ in Eq.~\eqref{eq:time-align} maps the spectral noise expansion output $\mathbf{x}^{s_{i+1}}_{t_i}$ to a valid flow-matching state $\tilde{\mathbf{x}}^{s_{i+1}}_{\tilde{t}_i}$ at the new resolution scale $s_{i+1} > s_i$. Throughout this subsection we write $r := s_{i+1}/s_i > 1$ for the resolution ratio; the final formulas are stated in both the $r$ and $s_{i+1}/s_i$ forms to tie directly back to Eqs.~\eqref{eq:scale-align}--\eqref{eq:time-align} in Sec.~\ref{sec:method:noise-injection}. The derivations below assume the ideal setting in which the network outputs match the ground-truth flow targets.

\paragraph{1. Setup.}
The spectral noise expansion produces an expanded state $\mathbf{x}^{s_{i+1}}_{t_i}$ in the frequency domain of the expanded resolution grid. For a grid with $N_{s_{i+1}}$ points, the coefficients are:
$$
\left(x^{s_{i+1}}_{t_i}\right)^{(\omega)} =
\begin{cases}
(1-t_i)\left(x_{0}^{s_i}\right)^{(\omega)} + t_i\epsilon^{(\omega)}, & \omega \in \Omega_{s_i}, \\[2pt]
t_i\epsilon'^{(\omega)}, & \omega \in \Omega_{s_{i+1}} \setminus \Omega_{s_i},
\end{cases}
$$
where $\epsilon^{(\omega)}, \epsilon'^{(\omega)} \sim \mathcal{N}(0, 1)$ are independent standard Gaussian variables. Under the orthonormal identification of spectra across resolutions, the clean-signal coefficient at a shared frequency $\omega \in \Omega_{s_i}$ relates to its scale-$s_{i+1}$ counterpart by the resolution ratio $r$:
$$
\left(x_{0}^{s_{i+1}}\right)^{(\omega)} = r \cdot \left(x_{0}^{s_i}\right)^{(\omega)} = (s_{i+1}/s_i) \cdot \left(x_{0}^{s_i}\right)^{(\omega)}, \quad \omega \in \Omega_{s_i}.
$$
For the newly exposed high-frequency slots $\omega \in \Omega_{s_{i+1}} \setminus \Omega_{s_i}$, the clean coefficients are treated as noise-dominated at timestep $t_i$ per Proposition~1, so spectral noise expansion initializes them with the correct noise level.

\paragraph{2. Timestep alignment.}
We seek a scaling factor $\kappa_i$ and an aligned timestep $\tilde{t}_i \in (0, 1)$ such that, for every $\omega \in \Omega_{s_{i+1}}$, the transformed state satisfies:
$$
\left(\tilde{x}^{s_{i+1}}_{\tilde{t}_i}\right)^{(\omega)} = \kappa_i\left(x^{s_{i+1}}_{t_i}\right)^{(\omega)} = (1-\tilde{t}_i)\left(x_{0}^{s_{i+1}}\right)^{(\omega)} + \tilde{t}_i\tilde\epsilon^{(\omega)}, \quad \tilde\epsilon^{(\omega)} \sim \mathcal{N}(0,1) \text{ i.i.d.}
$$
By matching the noise coefficients across all bands and the signal coefficients on the shared band $\Omega_{s_i}$, we derive two scalar conditions:
$$
\kappa_i \cdot t_i = \tilde{t}_i, \quad \kappa_i(1-t_i) = r(1-\tilde{t}_i).
$$
Substituting $\kappa_i = \tilde{t}_i/t_i$ from the noise match into the signal match equation gives:
$$
\frac{\tilde{t}_i}{t_i}(1-t_i) = r(1-\tilde{t}_i)
$$
Solving for $\tilde{t}_i$ and re-expressing the result in $s_{i+1}/s_i$ notation gives:
$$
\tilde{t}_i = \frac{r \cdot t_i}{1 + (r - 1)t_i} = \frac{(s_{i+1}/s_i) \cdot t_i}{1 + ((s_{i+1}/s_i) - 1)t_i}.
$$
Consequently, the scaling factor $\kappa_i$ is:
$$
\kappa_i = \frac{r}{1 + (r - 1)t_i} = \frac{s_{i+1}/s_i}{1 + ((s_{i+1}/s_i) - 1)t_i}.
$$
Substituting these back into the setup confirms that $\left(\tilde{x}^{s_{i+1}}_{\tilde{t}_i}\right)^{(\omega)} = (1-\tilde{t}_i)\left(x_{0}^{s_{i+1}}\right)^{(\omega)} + \tilde{t}_i\tilde\epsilon^{(\omega)}$ holds for all $\omega \in \Omega_{s_{i+1}}$. This proves that the resulting scaled and timestep-aligned state $\tilde{\mathbf{x}}^{s_{i+1}}_{\tilde{t}_i}$ is a valid flow-matching state at resolution $s_{i+1}$ and timestep $\tilde{t}_i$. $\square$

\subsection{Proof of Proposition 1: Per-frequency $\delta$-optimal Activation Time}
\label{app:proof-prop1}

We first establish a lemma giving the Bayes-optimal velocity error in closed form under a simplified modelling assumption $x_0^{(\omega)} \sim \mathcal{N}(0, P_\omega)$. In practice, we measure the power spectrum on centered (i.e. zero-mean) inputs $\mathbf{x}_0$ to get the power spectrum $P_\omega$, and this modelling assumption serves as a second moment approximation that allows the analytical derivation of our $\delta$-optimal resolution schedule. The proof of Proposition~1 then follows by solving the $\delta$-bound on the Bayes-optimal velocity error for $t$ to obtain the activation time $t_\omega$ in Eq.~\eqref{eq:activation}.

\begin{lemma}[Bayes-optimal velocity error] Let $v^{*(\omega)}$ be the per-frequency component of the Bayes-optimal velocity predictor $\mathbf{v}^{*}$ defined in Sec.~\ref{sec:method:prelim:basics}:

\begin{equation}
  \mathbf{v}^*(\mathbf{x}_t, t) \;=\; \mathbb{E}\!\left[\,\boldsymbol{\epsilon} - \mathbf{x}_0 \,\big|\, \mathbf{x}_t\,\right].
  \label{eq:optimal-v}
\end{equation}
Under the modelling assumption $x_0^{(\omega)} \sim \mathcal{N}(0, P_\omega)$, $v^{*(\omega)}$ deviates from the noise-only prediction $\epsilon^{(\omega)}$ in expectation by
\begin{equation}
  \mathbb{E}\!\left[\bigl|v^{*(\omega)}(\mathbf{x}_t, t) - \epsilon^{(\omega)}\bigr|^2\right]
  \;=\; \frac{\mathrm{SNR}_\omega(t)\,(1 + P_\omega)}{1 + \mathrm{SNR}_\omega(t)}.
  \label{eq:velocity-error-optimal}
\end{equation}
\end{lemma}


%\paragraph{Remark.} Real-data latents are not perfectly Gaussian, so Eq.~\eqref{eq:velocity-error-optimal} plays the role of an analytically tractable reference curve for our resolution schedule calculations. In practice, classifier-free guidance (CFG) from the text prompt during the generation process can also allow lower frequencies to be generated faster. We empirically measure this error in Sec.~\ref{sec:motivation}. \hx{TODO!}

\begin{proof}
\leavevmode

\paragraph{1. Assumptions.} The proof assumes a linear-Gaussian model, where:
\begin{enumerate}
  \item[(A1)] we model the clean data spectrum as zero-mean Gaussian $x_0^{(\omega)} \sim \mathcal{N}(0,\, P_\omega)$ with second moment $P_\omega > 0$;
  \item[(A2)] we assume $\epsilon^{(\omega)} \sim \mathcal{N}(0, 1)$, by orthonormality of $\Phi$ applied to $\boldsymbol{\epsilon} \sim \mathcal{N}(0, I)$;
  \item[(A3)] we assume $x_0^{(\omega)} \perp \epsilon^{(\omega)}$, i.e. the signal and noise coefficients are independent.
\end{enumerate}
(A1) is the simplified modelling assumption; (A2) and (A3) follow directly from the flow-matching forward process. Under (A1), the triple $(x_0^{(\omega)},\, \epsilon^{(\omega)},\, x_t^{(\omega)})$ is jointly Gaussian, so conditional expectations are linear and closed-form.

\paragraph{2. Second-moment quantities.} From the per-frequency forward process in Eq.~\eqref{eq:forward-mode} and (A1)--(A3),
\begin{align}
  \mathbb{E}[x_t^{(\omega)}]                     &= 0, \qquad
  \mathrm{Var}(x_t^{(\omega)}) \;=\; (1-t)^2 P_\omega + t^2,  \\
  \mathrm{Cov}(x_t^{(\omega)},\, x_0^{(\omega)})      &= (1-t)\,P_\omega, \qquad
  \mathrm{Cov}(x_t^{(\omega)},\, \epsilon^{(\omega)}) \;=\; t.
\end{align}

\paragraph{3. Bayes-optimal velocity as a linear predictor.}
Under (A1)--(A3), the triplet $(x_0^{(\omega)},\, \epsilon^{(\omega)},\, x_t^{(\omega)})$ is jointly Gaussian with zero mean, since $x_t^{(\omega)} = (1-t) x_0^{(\omega)} + t\,\epsilon^{(\omega)}$ is a linear combination of two independent zero-mean Gaussians. For any pair of zero-mean jointly Gaussian random variables $(X, Y)$ with $\mathrm{Var}(Y) > 0$, the conditional expectation is a linear function of $Y$:
\begin{equation}
  \mathbb{E}[X \,|\, Y] \;=\; \frac{\mathrm{Cov}(X, Y)}{\mathrm{Var}(Y)}\,Y.
  \label{eq:gauss-cond}
\end{equation}

We apply \eqref{eq:gauss-cond} twice, using the second moments computed above. First, with $X = x_0^{(\omega)}$ and $Y = x_t^{(\omega)}$,
\begin{equation}
  \mathbb{E}\bigl[x_0^{(\omega)} \,\big|\, x_t^{(\omega)}\bigr]
  \;=\; \frac{(1-t)\,P_\omega}{(1-t)^2 P_\omega + t^2}\, x_t^{(\omega)}.
  \label{eq:cond-x0}
\end{equation}
Second, with $X = \epsilon^{(\omega)}$ and $Y = x_t^{(\omega)}$,
\begin{equation}
  \mathbb{E}\bigl[\epsilon^{(\omega)} \,\big|\, x_t^{(\omega)}\bigr]
  \;=\; \frac{t}{(1-t)^2 P_\omega + t^2}\, x_t^{(\omega)}.
  \label{eq:cond-eps}
\end{equation}

The Bayes-optimal velocity at frequency $\omega$ is $v^{*(\omega)} = \mathbb{E}\bigl[\epsilon^{(\omega)} - x_0^{(\omega)} \,\big|\, \mathbf{x}_t\bigr]$, the $\phi_\omega$-projection of the full-input optimal velocity in Eq.~\eqref{eq:optimal-v}. Under the modelling assumption that $\mathbf{x}_0$ has diagonal covariance in basis $\Phi$ (i.e., (A1) holds independently for every $\omega$), the joint Gaussian $(\mathbf{x}_0, \boldsymbol{\epsilon}, \mathbf{x}_t)$ factorizes across $\Omega$, so the posterior collapses to its single-coefficient marginal:
\begin{equation}
  \mathbb{E}\bigl[\epsilon^{(\omega)} - x_0^{(\omega)} \,\big|\, \mathbf{x}_t\bigr]
  \;=\; \mathbb{E}\bigl[\epsilon^{(\omega)} - x_0^{(\omega)} \,\big|\, x_t^{(\omega)}\bigr].
  \label{eq:full-to-scalar}
\end{equation}
Subtracting \eqref{eq:cond-x0} from \eqref{eq:cond-eps} and combining with \eqref{eq:full-to-scalar},
\begin{equation}
  v^{*(\omega)} \;=\; A \cdot x_t^{(\omega)},
  \qquad
  A \;:=\; \frac{t - (1-t)\,P_\omega}{(1-t)^2 P_\omega + t^2}.
  \label{eq:vstar-linear}
\end{equation}

\paragraph{4. Substitute and simplify.} Plugging $x_t^{(\omega)} = (1-t)\, x_0^{(\omega)} + t\,\epsilon^{(\omega)}$ into \eqref{eq:vstar-linear} and collecting terms in $x_0^{(\omega)}$ and $\epsilon^{(\omega)}$,
\begin{equation}
  v^{*(\omega)} - \epsilon^{(\omega)}
  \;=\; A\bigl[(1-t)\,x_0^{(\omega)} + t\,\epsilon^{(\omega)}\bigr] - \epsilon^{(\omega)}
  \;=\; A(1-t)\,x_0^{(\omega)} \,+\, (At - 1)\,\epsilon^{(\omega)}.
\end{equation}
Squaring and taking expectation, using (A1)--(A3) (so $\mathbb{E}[(x_0^{(\omega)})^2] = P_\omega$, $\mathbb{E}[(\epsilon^{(\omega)})^2] = 1$, and $\mathbb{E}[x_0^{(\omega)}\,\epsilon^{(\omega)}] = 0$), the cross term vanishes:
\begin{equation}
  \mathbb{E}\!\left[\bigl|v^{*(\omega)} - \epsilon^{(\omega)}\bigr|^2\right]
  \;=\; A^2 (1-t)^2\,P_\omega \,+\, (At - 1)^2.
\end{equation}
Expanding $(At-1)^2 = A^2 t^2 - 2At + 1$ and grouping the $A^2$ terms,
\begin{equation}
  \mathbb{E}\!\left[\bigl|v^{*(\omega)} - \epsilon^{(\omega)}\bigr|^2\right]
  \;=\; A^2\!\left[(1-t)^2 P_\omega + t^2\right] - 2At + 1.
  \label{eq:verror-expanded}
\end{equation}
Denote $D := (1-t)^2 P_\omega + t^2 = \mathrm{Var}(x_t^{(\omega)})$; from Eq.~\eqref{eq:vstar-linear},
\begin{equation}
  A \;=\; \frac{t - (1-t)\, P_\omega}{D}.
  \label{eq:A-in-D}
\end{equation}
Substituting into Eq.~\eqref{eq:verror-expanded} term by term,
\begin{align}
  A^2\!\left[(1-t)^2 P_\omega + t^2\right] \;&=\; A^2 D \;=\; \frac{(t - (1-t)\, P_\omega)^2}{D}, \\
  2At \;&=\; \frac{2t\,(t - (1-t)\, P_\omega)}{D},
\end{align}
so the right-hand side of Eq.~\eqref{eq:verror-expanded} combines into a single fraction over $D$:
\begin{equation}
  \mathbb{E}\!\left[\bigl|v^{*(\omega)} - \epsilon^{(\omega)}\bigr|^2\right]
  \;=\; \frac{(t - (1-t)\, P_\omega)^2 \,-\, 2t\,(t - (1-t)\, P_\omega) \,+\, D}{D}.
  \label{eq:verror-singlefrac}
\end{equation}
Expanding $(t - (1-t)\, P_\omega)^2 = t^2 - 2t(1-t) P_\omega + (1-t)^2 P_\omega^2$, the numerator of Eq.~\eqref{eq:verror-singlefrac} simplifies as
\begin{align}
  & \bigl[t^2 - 2t(1-t) P_\omega + (1-t)^2 P_\omega^2\bigr]
    - \bigl[2t^2 - 2t(1-t) P_\omega\bigr]
    + \bigl[(1-t)^2 P_\omega + t^2\bigr] \nonumber\\
  ={}& (1-t)^2 P_\omega^2 + (1-t)^2 P_\omega
    \;=\; (1-t)^2\,P_\omega\,(1 + P_\omega).
\end{align}
Substituting back into Eq.~\eqref{eq:verror-singlefrac},
\begin{equation}
  \mathbb{E}\!\left[\bigl|v^{*(\omega)} - \epsilon^{(\omega)}\bigr|^2\right]
  \;=\; \frac{(1-t)^2\,P_\omega\,(1 + P_\omega)}{(1-t)^2 P_\omega + t^2}.
\end{equation}

\paragraph{5. Rewrite via $\mathrm{SNR}_\omega(t)$.} Using $\mathrm{SNR}_\omega(t) = (1-t)^2 P_\omega / t^2$ from Eq.~\eqref{eq:snr}, the denominator is $(1-t)^2 P_\omega + t^2 = t^2(1 + \mathrm{SNR}_\omega(t))$ and the numerator equals $t^2\,\mathrm{SNR}_\omega(t)\,(1 + P_\omega)$. The $t^2$ cancels, recovering the right-hand side of Eq.~\eqref{eq:velocity-error-optimal} and completing the proof of the lemma.

\end{proof}

\begin{proposition}[Per-frequency $\delta$-optimal activation time] Denote the per-frequency Bayes-optimal velocity predictor by $v^{*(\omega)}(\mathbf{x}_t, t)$, whose spatial-domain counterpart is defined in Eq.~\eqref{eq:optimal-v}. Under the simplified modelling assumption $x_0^{(\omega)} \sim \mathcal{N}(0, P_\omega)$ with $P_\omega > 0$, we have
\begin{equation}
  \mathbb{E}\!\left[\bigl|v^{*(\omega)}(\mathbf{x}_t, t) - \epsilon^{(\omega)}\bigr|^2\right] \;\le\; \delta
  \label{eq:noise-dominated-appendix}
\end{equation}
for all $t \ge t_\omega$ and $\delta\in(0, 1)$, where
\begin{equation}
  t_\omega \;:=\; \frac{1}{1 + \sqrt{\dfrac{\delta}{P_\omega\,(1 + P_\omega - \delta)}}}.
  \label{eq:activation-appendix}
\end{equation}
\end{proposition}

\begin{proof}
\leavevmode

\paragraph{Solving the $\delta$-bound for $t$.} Substituting the closed-form error into the tolerance bound $\mathbb{E}[|v^{*(\omega)} - \epsilon^{(\omega)}|^2] \le \delta$ of Proposition~1 yields
\begin{equation}
  \frac{\mathrm{SNR}_\omega(t)\,(1 + P_\omega)}{1 + \mathrm{SNR}_\omega(t)} \;\le\; \delta.
\end{equation}
Clearing the (positive) denominator and collecting terms in $\mathrm{SNR}_\omega(t)$,
\begin{equation}
  \mathrm{SNR}_\omega(t)\,(1 + P_\omega - \delta) \;\le\; \delta,
  \qquad \text{i.e.,} \qquad
  \mathrm{SNR}_\omega(t) \;\le\; \frac{\delta}{1 + P_\omega - \delta},
\end{equation}
where the rearrangement is valid because $\delta < 1 \le 1 + P_\omega$ keeps $1 + P_\omega - \delta > 0$. Substituting $\mathrm{SNR}_\omega(t) = (1-t)^2 P_\omega / t^2$ and taking the positive square root (all quantities positive for $t \in (0, 1)$),
\begin{equation}
  \frac{1-t}{t} \;\le\; \sqrt{\frac{\delta}{P_\omega\,(1 + P_\omega - \delta)}}.
\end{equation}
Adding $1$ to both sides gives $1/t \le 1 + \sqrt{\delta / (P_\omega(1 + P_\omega - \delta))}$, which inverts to
\begin{equation}
  t \;\ge\; \frac{1}{1 + \sqrt{\dfrac{\delta}{P_\omega\,(1 + P_\omega - \delta)}}} \;=\; t_\omega,
\end{equation}
recovering Eq.~\eqref{eq:activation-appendix}.
\end{proof}


\subsection{Proof of Proposition 2: Per-resolution $\delta$-optimal Transition Time}
\label{app:proof-prop2}

\begin{proposition}[Per-resolution $\delta$-optimal transition time] Under the setting of Proposition~1, and assuming $P_\omega$ is monotonically decreasing in $|\omega|$ based on the power-law decay in Eq.~\eqref{eq:power-law}, for any resolution scale $s_i, s_{i+1} \in (0, 1]$, the optimal transition time from scale $s_i$ up to $s_{i+1}$ is
\begin{equation}
  t_i^* \;:=\; \min_{\omega \in \Omega_{s_i}}\, t_\omega \;=\; t_{\omega=s_i \cdot \omega_{\max}(H,W)},
  \label{eq:stage-transition-appendix}
\end{equation}
where $\Omega_{s_i}$ is the set of frequencies representable on the $(s_iH, s_iW)$ grid.
\end{proposition}

\begin{proof}
\leavevmode

We show that the activation time $t_\omega$ of Eq.~\eqref{eq:activation} is strictly decreasing in $|\omega|$ under the assumptions of the proposition; the claim then follows immediately.

\paragraph{1. $t_\omega$ is strictly increasing in $P_\omega$.}
For $\delta \in (0, 1)$ and $P_\omega > 0$, the quadratic $P_\omega\,(1 + P_\omega - \delta) = P_\omega^2 + (1-\delta)\, P_\omega$ has derivative $2 P_\omega + (1-\delta) > 0$, hence is strictly increasing in $P_\omega$. Therefore the radicand $\delta / \bigl(P_\omega(1 + P_\omega - \delta)\bigr)$ in Eq.~\eqref{eq:activation} is strictly decreasing in $P_\omega$, and so is its square root. Hence
\begin{equation}
  t_\omega \;=\; \frac{1}{1 + \sqrt{\dfrac{\delta}{P_\omega\,(1 + P_\omega - \delta)}}}
\end{equation}
is strictly increasing in $P_\omega$.

\paragraph{2. $P_\omega$ is strictly decreasing in $|\omega|$.}
The power-law assumption in Eq.~\eqref{eq:power-law} gives $P_\omega \propto |\omega|^{-\beta}$ with $\beta > 0$. Thus $P_\omega$ is strictly decreasing in $|\omega|$.

\paragraph{3. Minimum over $\Omega_s$.}
Combining the above observations, $t_\omega$ is strictly decreasing in $|\omega|$. The set $\Omega_s = \{\omega \in \Omega : |\omega| \le s \cdot \omega_{\max}(H, W)\}$ consists exactly of the frequencies with radius at most $s \cdot \omega_{\max}(H, W)$, so
\begin{equation}
  \min_{\omega \in \Omega_s}\, t_\omega \;=\; t_{\,s \cdot \omega_{\max}(H, W)},
\end{equation}
which is Eq.~\eqref{eq:stage-transition}.
\end{proof}

\clearpage
\section{Empirical Validation of the Per-frequency Activation Criterion}
\label{sec:motivation}

We empirically validate Proposition 1 with a simple \textit{spectral noise passthrough} experiment.
While maintaining full-resolution processing at all denoising steps, for each frequency $\omega \in \Omega$, if the current timestep $t$ is before its activation time $t_\omega$ (i.e., $t > t_\omega$), we replace that frequency component with $t \cdot T_\Phi(\boldsymbol{\epsilon})^{(\omega)}$, where $\boldsymbol{\epsilon} \sim \mathcal{N}(0, I)$ is the initial noise at $t = 1$. We use the Discrete Cosine Transform (DCT) as $T_\Phi$ in this experiment and sweep over multiple $\delta$ values. As described in Sec.~\ref{sec:method:schedule:signal}, a larger $\delta$ corresponds to transitioning to the high-resolution stage later in denoising with more speedup, while a smaller $\delta$ indicates transitioning to high-resolution earlier in denoising with less speedup.

From Fig. \ref{fig:spectral-passthrough}, we observe that at smaller $\delta \in [0.0001, 0.001]$, there is almost no observable difference compared to native full-resolution generation. As $\delta$ increases, we see a clear transition to blurry and lower-quality images around $\delta=0.05$, and larger values further degrade generation quality. This finding provides empirical support for the spectral autoregression property of diffusion models and Proposition 1, which serves as the foundation of our Spectral Progressive Diffusion framework and the $\delta$-optimal resolution schedule. It also implies that the Gaussian modelling assumption in Proposition 1 serves as a reasonable approximation for subsequent analytical resolution schedule derivations.

\begin{figure}[h!]
    \centering
    \includegraphics[width=\linewidth]{supp_figures/spectral_passthrough.pdf}
    \caption{\textbf{Spectral noise passthrough experiment.} At smaller $\delta \in [0.0001, 0.001]$, there is almost no observable difference compared to native full-resolution generation. As larger $\delta$ values cause high-frequency replacement to persist later in the denoising trajectory, we observe increasingly blurry and distorted results (i.e., ghosting artifacts and ``CHOOLBUS'' instead of ``SCHOOLBUS''). }
    \label{fig:spectral-passthrough}
\end{figure}


\clearpage
\section{Additional Details on Power Spectrum Measurement}
\label{sec:power_spec}
To remain consistent with the zero-mean Gaussian modelling assumption in Sec.~\ref{sec:method:schedule:signal}, we center each data sample before measuring the radially averaged power spectrum of the VAE latent space of FLUX.1-dev~\citep{black2024flux} and Z-Image~\citep{zimage} for image generation, and the VAE latent space of WAN 2.1~\citep{wan} for video generation.
We fit the power-law model $P_\omega \approx A \cdot |\omega|^{-\beta}$ and report the fitted $A$ and $\beta$ values for image latents, video latents, and image pixels in Table~\ref{tab:power-spec-fit}.
These fitted values instantiate the per-frequency signal power $P_\omega$ in the activation-time and transition-time formulas in Sec.~\ref{sec:method:schedule:signal}, which we use to compute the $\delta$-optimal resolution schedule for each model family.

For image latents, we run the FLUX.1-dev VAE on 100K images from the Aesthetics-Train-V2~\citep{zhang2025diffusion4k} dataset, center-cropped and resized to $1024^2$ resolution. We then average power over channels and radial frequency bins before fitting the spectrum. For video latents, we use the same procedure with the WAN 2.1 VAE on 10K videos from the VChitect-T2V-Dataverse~\citep{fan2025vchitect} dataset, averaging over channels, latent frames, and radial frequency bins to estimate the spatial power spectrum. We further measure the pixel-space power spectrum on images from the Aesthetics-Train-V2 dataset, center-cropped and resized to $512^2$, the native resolution supported by the PixelGen model.

\begin{table}[h!]
  \centering
  \small
  \caption{\textbf{Power-law fits to measured power spectra.} We fit $P_\omega \approx A \cdot |\omega|^{-\beta}$ on centered samples for image latents, video latents, and image pixels.}
  \label{tab:power-spec-fit}
  \resizebox{\linewidth}{!}{
  \begin{tabular}{lcccc}
    \toprule
    Spectrum & $\beta$ & $A$ & $R^2$ & Samples \\
    \midrule
    FLUX.1-dev (Z-Image) latent & 1.9155 & 203.62 & 0.978 & 100K images \\
    WAN latent & 2.4227 & 219.48 & 0.999 & 10K videos \\
    Pixel $512^2$ & 2.4493 & 3745.46 & 1.000 & 100K images \\
    \bottomrule
  \end{tabular}
  }
\end{table}


\clearpage
\section{Additional Details and Results on Image Generation}
\label{sec:image_generation}


\subsection{Training-free Acceleration of Latent-space Image Generation}
\label{subsec:image_latent}

For training-free inference acceleration of image generation, we use NVIDIA A100 GPUs with 80~GB of VRAM to match RALU~\citep{ralu}'s evaluation protocol. For training-free inference acceleration of video generation, as well as image editing and ablation studies, we use NVIDIA H100 GPUs with 80~GB of VRAM. For training-free acceleration of latent-space image generation, we use FLUX.1-dev~\citep{black2024flux} as the base model and follow the evaluation protocol of RALU~\citep{ralu}. Table~\ref{tab:flux-extended} extends Table~\ref{tab:flux-main} with additional training-free acceleration baselines from RALU. We organize the comparison into three wall-clock speedup tiers:

\begin{itemize}
    \item \textbf{Small speedup}: the 50-step native full-resolution baseline, compared against RALU (50 steps) and our method (40 steps), yielding approximately $1.7\times$ speedup;
    \item \textbf{Medium speedup}: the 10-step native full-resolution baseline, used as an iso-compute reference for RALU (15 steps) and our method (12 steps), yielding approximately $5\times$ speedup; and
    \item \textbf{Large speedup}: the 7-step native full-resolution baseline, used as an iso-compute reference for RALU (10 steps) and our method (10 steps), yielding approximately $7\times$ speedup.
\end{itemize}
All speedups are calculated relative to the wall-clock latency of the 50-step FLUX.1-dev baseline on a single A100 GPU, following RALU~\citep{ralu}. As in Table~\ref{tab:flux-main}, we evaluate image quality using CLIP-IQA~\citep{wang2023clipiqa} and NIQE~\citep{mittal2013niqe}, prompt alignment using T2I-CompBench~\citep{t2icompbench} and GenEval~\citep{geneval}, and overall quality using ImageReward~\citep{imagereward}.
Following RALU's evaluation protocol, ImageReward, CLIP-IQA, and NIQE are averaged over 5,000 images generated from MS-COCO validation prompts. For compositional alignment, GenEval is evaluated with four random seeds per prompt on its object-focused prompt suite, and T2I-CompBench is evaluated on its spatial, non-spatial, and complex prompt subsets with four random seeds per prompt. Unless otherwise specified, all metrics are computed using their official or default evaluation settings.

We keep the error threshold fixed at $\delta=0.01$, use $T_\Phi = \text{DCT}$, and compare $S=2$ and $S=3$ resolution-stage settings. To align with the multi-resolution training distribution of FLUX.1-dev, we use latent dimensions $64^2$ ($s_1 = 0.5$) and $128^2$ $(s_2 = 1.0)$ (pixel dimensions $512^2$ and $1024^2$) for $S=2$, and latent dimensions $32^2$ $(s_1 = 0.25)$, $64^2$ $(s_2 = 0.5)$, and $128^2$ ($s_3 = 1.0$) (pixel dimensions $256^2$, $512^2$, and $1024^2$) for $S=3$.

By combining spectral transformations for progressive-resolution generation with our derived resolution schedule, our method achieves a better speed-quality tradeoff than the reduced-step FLUX.1-dev baselines and existing temporal and spatial diffusion acceleration methods, including progressive-resolution methods such as RALU~\citep{ralu}. In the largest speedup tier, our $S=3$ setting reaches $7.09\times$ wall-clock speedup and $7.36\times$ FLOPs speedup while preserving competitive image quality. We further provide extended qualitative baseline comparisons in Figures~\ref{fig:flux-1}--\ref{fig:flux-3}. FLUX.1-dev with reduced steps degrades image quality and exhibits over-saturation artifacts, while RALU introduces noticeable noise artifacts in reduced-step settings.


\begin{table*}[t!]
    \centering
    \caption{\textbf{Extended quantitative comparisons on training-free latent-space image generation.} We evaluate FLUX.1-dev at $1024^2$ resolution and group methods into the same $1\times$, $5\times$, and $7\times$ wall-clock speedup tiers as Table~\ref{tab:flux-main}. Rows shared with Table~\ref{tab:flux-main} use the same values. Speedup is calculated relative to the 50-step FLUX.1-dev baseline wall-clock runtime on a single A100 GPU. \textcolor{blue}{\textbf{T}} and \textcolor{red}{\textbf{S}} denote temporal and spatial acceleration methods, respectively. Baseline metrics are copied from RALU~\citep{ralu} under the same evaluation protocol.}
    \label{tab:flux-extended}
    \resizebox{\textwidth}{!}{
    \begin{scriptsize}
    \begin{tabular}{cccccccccc}
         \toprule
        \multirow{2}{*}{Method} & \multirow{2}{*}{Accel.} & \multirow{2}{*}{Latency (s) $\downarrow$} & \multirow{2}{*}{Speedup (s) $\uparrow$} & \multirow{2}{*}{TFLOPs $\downarrow$} & Overall & \multicolumn{2}{c}{Image quality} & \multicolumn{2}{c}{Text alignment} \\
        \cmidrule(lr){6-10}
        & & & & & ImageReward $\uparrow$ & CLIP-IQA $\uparrow$ & NIQE $\downarrow$ & T2I-Comp. $\uparrow$ & GenEval $\uparrow$  \\
        \midrule
        FLUX (50 steps) & - & 25.1 & 1.00$\times$ & 2991.01 & \textbf{1.095} & 0.707 & 6.75 & \textbf{0.634} & \textbf{0.698}  \\
        RALU~\citep{ralu}   & \textcolor{red}{\textbf{S}}  & 15.85 & 1.58$\times$ & \underline{1749.94} & 1.028 & \underline{0.712} & \textbf{6.07} & 0.613 & 0.648 \\
        \textbf{Ours ($S=2$)}     & \textcolor{red}{\textbf{S}}       & \underline{15.16} & \underline{1.66}$\times$ & 1755.22 & \underline{1.049} & \textbf{0.719} & 6.43 & \underline{0.617} & \underline{0.654} \\
        \textbf{Ours ($S=3$)}     & \textcolor{red}{\textbf{S}}       & \textbf{14.54} & \textbf{1.73}$\times$ & \textbf{1672.04} & 1.015 & 0.711 & \underline{6.33} & 0.593 & 0.640 \\
        \cmidrule(lr){1-10}
        FLUX (10 steps) & \textcolor{blue}{\textbf{T}} & 5.18 & 4.84$\times$ & 610.02 & 0.981 & 0.679 & 6.93 & 0.618 & 0.647 \\
        $\Delta$-DiT~\citep{deltadit}      & \textcolor{blue}{\textbf{T}} & 7.42 & 3.38$\times$ & 772.10 & 0.102 & 0.487 & 9.60 & 0.306 & 0.397 \\
        ToCa~\citep{toca}                  & \textcolor{blue}{\textbf{T}} & 15.5 & 1.62$\times$ & 601.12 & -1.827 & 0.253 & 10.6 & 0.259 & 0.137 \\
        TeaCache~\citep{teacache}          & \textcolor{blue}{\textbf{T}} & 5.23 & 4.80$\times$ & 610.59 & 0.944 & 0.665 & 7.92 & 0.620 & 0.647 \\
        TaylorSeer~\citep{taylorseer}      & \textcolor{blue}{\textbf{T}} & 9.34 & 2.69$\times$ & 556.72 & 0.972 & 0.684 & 6.77 & 0.594 & 0.619 \\
        Bottleneck~\citep{bottleneck}      & \textcolor{red}{\textbf{S}}  & 5.37 & 4.67$\times$ & 571.23 & 0.889 & 0.661 & 9.16 & 0.620 & \textbf{0.687} \\
        RALU~\citep{ralu}                  & \textcolor{red}{\textbf{S}}  & 5.04 & 4.98$\times$ & 540.47 & 1.022 & \underline{0.700} & \textbf{6.43} & \textbf{0.626} & 0.652 \\
        \textbf{Ours ($S=2$)}     & \textcolor{red}{\textbf{S}}       & \underline{4.35} & \underline{5.77}$\times$ & \underline{500.34} & \textbf{1.059} & 0.696 & 6.69 & \underline{0.624} & \underline{0.655} \\
        \textbf{Ours ($S=3$)}     & \textcolor{red}{\textbf{S}}       & \textbf{4.12} & \textbf{6.09}$\times$ & \textbf{469.15} & \underline{1.042} & \textbf{0.701} & \underline{6.53} & 0.623 & 0.637 \\
        \cmidrule(lr){1-10}
        FLUX (7 steps) & \textcolor{blue}{\textbf{T}} & 3.79 & 6.62$\times$ & 431.45 & 0.920 & 0.660 & 8.25 & 0.594 & 0.583 \\
        TeaCache~\citep{teacache}          & \textcolor{blue}{\textbf{T}} & 4.21 & 5.96$\times$ & 431.83 & 0.733 & 0.623 & 13.7 & 0.599 & 0.594 \\
        TaylorSeer~\citep{taylorseer}      & \textcolor{blue}{\textbf{T}} & 7.00 & 3.59$\times$ & 431.74 & 0.660 & 0.646 & 9.43 & 0.514 & 0.446 \\
        Bottleneck~\citep{bottleneck}      & \textcolor{red}{\textbf{S}}  & 3.78 & 6.64$\times$ & 431.52 & 0.792 & 0.631 & 8.71 & 0.605 & \underline{0.672} \\
        RALU~\citep{ralu}                  & \textcolor{red}{\textbf{S}}  & 3.75 & 6.69$\times$ & \underline{426.01} & 0.999 & 0.681 & 6.87 & \textbf{0.633} & \textbf{0.682} \\
        \textbf{Ours ($S=2$)}     & \textcolor{red}{\textbf{S}}       & \underline{3.70} & \underline{6.78}$\times$ & 427.03 & \textbf{1.039} & \underline{0.689} & \underline{6.78} & 0.620 & 0.667 \\
        \textbf{Ours ($S=3$)}     & \textcolor{red}{\textbf{S}}       & \textbf{3.54} & \textbf{7.09}$\times$ & \textbf{406.24} & \underline{1.015} & \textbf{0.694} & \textbf{5.99} & \underline{0.627} & 0.637 \\
        \bottomrule
    \end{tabular}
    \end{scriptsize}}
\end{table*}


\subsection{Spectral-transformation-based Fine-tuning for Latent- and Pixel-Space Image Generation}
\label{subsec:image_pixel}
\paragraph{Additional latent-space fine-tuning results.} Since FLUX.1-dev is a guidance-distilled model and its base model is not open-sourced, we additionally conduct latent-space image generation fine-tuning experiments on Z-Image~\citep{zimage}, which employs the same VAE architecture as FLUX.1-dev and hence the same $\delta$-optimal resolution schedules. We adopt Low-Rank Adaptation (LoRA) of rank 32 and train for 2000 steps following the spectral-transformation-based fine-tuning procedure in Sec.~\ref{sec:method:lora}. For fair comparison with the full-resolution and training-free acceleration baselines, we fine-tune on synthetic images generated by full-resolution Z-Image with 50 denoising steps using a set of 5K MS-COCO prompts that is disjoint from the 5K MS-COCO validation prompts used for evaluation. Training is conducted on our internal cluster of 8 NVIDIA H100 GPUs, each with 80~GB of VRAM, at an effective batch size of 8 and takes approximately 2 hours. We keep the optimizer and all other LoRA fine-tuning hyperparameters at the default settings of the \texttt{DiffSynth Studio} library, including a constant learning rate of 1e-4. Fine-tuning compute scales linearly with the number of sampled training images, training iterations, and batch size. In our experiments we fix the synthetic fine-tuning set to 5K prompts and train LoRA adapters for 2000 iterations, so increasing the fine-tuning dataset size would only affect compute through additional sampled training examples and does not change the inference-time cost of our method.

All latent-image generation experiments are conducted at a native resolution of $1024^2$ pixels, which is the default configuration for both FLUX.1-dev and Z-Image. The progressive resolution settings for the $S = 2$ and $S = 3$ cases are the same as in the FLUX.1-dev training-free acceleration experiments in Sec.~\ref{subsec:image_latent}.

Further model fine-tuning allows a more aggressive error tolerance $\delta$, leading to an even higher speedup. For fine-tuning experiments, we additionally evaluate $\delta=0.05$, which is larger than the default $\delta=0.01$. As in Sec.~\ref{subsec:image_latent}, we organize the comparison into three wall-clock speedup tiers, all calculated relative to the 50-step Z-Image baseline on a single H100 GPU (we adopt H100 wall-clock measurements for Z-Image and all subsequent pixel-space image generation and video generation experiments):

\begin{itemize}
    \item \textbf{Small speedup}: the 50-step native full-resolution Z-Image baseline, compared against our training-free acceleration and LoRA fine-tuned variants at $S=2$ and $S=3$ ($\delta=0.01$, 50 steps), as well as our LoRA fine-tuned variant with the configuration ($\delta=0.05$, $S=2$, 50 steps), yielding $1.65\times$ to $2.01\times$ speedup;
    \item \textbf{Medium speedup}: the 10-step native full-resolution Z-Image baseline, used as an iso-compute reference for our training-free acceleration and LoRA fine-tuned variants at $S=2$ and $S=3$ with reduced denoising steps, yielding approximately $5\times$ speedup; and
    \item \textbf{Large speedup}: the 7-step native full-resolution Z-Image baseline, used as an iso-compute reference for our training-free acceleration and LoRA fine-tuned variants at $S=2$ and $S=3$ with further reduced denoising steps, yielding approximately $7\times$ to $8\times$ speedup.
\end{itemize}
For quantitative evaluation, we use the same metric evaluation settings as Sec.~\ref{subsec:image_latent}, following the RALU evaluation protocol, and report the numbers in Table~\ref{tab:zimage-supp}. We show extensive qualitative comparisons in Figures~\ref{fig:zimage-1}--\ref{fig:zimage-3}. Across all three speedup tiers, our training-free acceleration variants outperform the reduced-step Z-Image baselines, and our spectral-transformation-based LoRA fine-tuning further closes the gap to the 50-step full-resolution baseline while attaining up to $7.81\times$ wall-clock speedup. In particular, our LoRA fine-tuned model at $S=3$ outperforms the training-free acceleration variant at $S=2$ while improving efficiency, and our LoRA fine-tuned model at $S=2$ either with reduced denoising steps or at a larger $\delta=0.05$ attains image quality comparable to the training-free acceleration $S=2$, 50-step setting, providing additional efficiency at no quality cost.


\begin{table*}[t!]
    \centering
    \caption{\textbf{Quantitative comparisons on latent-space image generation with fine-tuning.} We evaluate Z-Image at $1024^2$ resolution and group methods into wall-clock speedup tiers. Speedup is calculated relative to the 50-step Z-Image baseline wall-clock runtime on a single H100 GPU. \textcolor{blue}{\textbf{T}} and \textcolor{red}{\textbf{S}} denote temporal and spatial acceleration methods, respectively. All metrics are evaluated using the same protocol as Table~\ref{tab:flux-main}.}
    \label{tab:zimage-supp}
    \resizebox{\textwidth}{!}{
    \begin{scriptsize}
    \begin{tabular}{cccccccccc}
         \toprule
        \multirow{2}{*}{Method} & \multirow{2}{*}{Accel.} & \multirow{2}{*}{Latency (s) $\downarrow$} & \multirow{2}{*}{Speedup (s) $\uparrow$} & \multirow{2}{*}{TFLOPs $\downarrow$} & Overall & \multicolumn{2}{c}{Image quality} & \multicolumn{2}{c}{Text alignment} \\
        \cmidrule(lr){6-10}
        & & & & & ImageReward $\uparrow$ & CLIP-IQA $\uparrow$ & NIQE $\downarrow$ & T2I-Comp. $\uparrow$ & GenEval $\uparrow$  \\
        \midrule
        Z-Image (50 steps) & - & 21.25 & 1.00$\times$ & 4941.23 & 0.965 & 0.700 & 5.41 & 0.731 & 0.745  \\
        \cmidrule(lr){1-10}
        Z-Image (32 steps) & \textcolor{blue}{\textbf{T}} & 13.59 & 1.56$\times$ & 3166.62 & \underline{0.957} & 0.697 & \textbf{5.44} & 0.686 & 0.725  \\
        \textbf{Ours (TF, $\delta=0.01$, $S=2$)}     & \textcolor{red}{\textbf{S}}       & 12.90 & 1.65$\times$ & 3132.03 & 0.904 & 0.688 & 5.87 & 0.658 & \underline{0.730} \\
        \textbf{Ours (TF, $\delta=0.01$, $S=3$)}     & \textcolor{red}{\textbf{S}}       & \underline{12.19} & \underline{1.74$\times$} & \underline{2871.09} & 0.875 & 0.690 & \underline{5.59} & 0.650 & 0.682 \\
        \textbf{Ours (LoRA, $\delta=0.01$, $S=2$)}     & \textcolor{red}{\textbf{S}}       & 12.87 & 1.65$\times$ & 3132.03 & \textbf{0.982} & \textbf{0.699} & 5.72 & \textbf{0.725} & \textbf{0.731} \\
        \textbf{Ours (LoRA, $\delta=0.01$, $S=3$)}     & \textcolor{red}{\textbf{S}}       & 12.23 & 1.74$\times$ & \underline{2871.09} & 0.954 & \underline{0.697} & 5.75 & \underline{0.717} & 0.728 \\
        \textbf{Ours (LoRA, $\delta=0.05$, $S=2$)}     & \textcolor{red}{\textbf{S}}       & \textbf{10.59} & \textbf{2.01$\times$} & \textbf{2436.19} & 0.919 & 0.684 & 6.12 & 0.661 & 0.718 \\
        \cmidrule(lr){1-10}
        Z-Image (10 steps) & \textcolor{blue}{\textbf{T}} & 4.26 & 4.99$\times$ & 997.68 & 0.851 & 0.659 & 5.95 & \underline{0.678} & \underline{0.705} \\
        \textbf{Ours (TF, $\delta=0.01$, $S=2$)}     & \textcolor{red}{\textbf{S}}       & \underline{4.22} & \underline{5.04$\times$} & \underline{962.95} & \underline{0.860} & \underline{0.668} & 6.17 & 0.677 & 0.693 \\
        \textbf{Ours (TF, $\delta=0.01$, $S=3$)}     & \textcolor{red}{\textbf{S}}       & \textbf{4.01} & \textbf{5.30$\times$} & \textbf{875.97} & 0.827 & 0.662 & \underline{5.78} & 0.676 & 0.655 \\
        \textbf{Ours (LoRA, $\delta=0.01$, $S=2$)}     & \textcolor{red}{\textbf{S}}       &  4.24 & 5.01$\times$ & \underline{962.95} & \textbf{0.923} & \textbf{0.683} & \textbf{5.58} & \textbf{0.706} & \textbf{0.738} \\
        \cmidrule(lr){1-10}
        Z-Image (7 steps) & \textcolor{blue}{\textbf{T}} & 3.56 & 5.97$\times$ & 701.92 & 0.763 & 0.631 & \textbf{6.14} & 0.648 & 0.667 \\
        \textbf{Ours (TF, $\delta=0.01$, $S=2$)}     & \textcolor{red}{\textbf{S}}       & 2.75 & 7.73$\times$ & \underline{678.77} & \underline{0.804} & \underline{0.641} & 6.46 & \underline{0.665} & \underline{0.680} \\
        \textbf{Ours (TF, $\delta=0.01$, $S=3$)}     & \textcolor{red}{\textbf{S}}       & \textbf{2.63} & \textbf{8.08$\times$} & \textbf{609.18} & 0.759 & 0.630 & \underline{6.16} & 0.658 & 0.645 \\
        \textbf{Ours (LoRA, $\delta=0.01$, $S=2$)}     & \textcolor{red}{\textbf{S}}       & \underline{2.72} & \underline{7.81$\times$} & \underline{678.77} & \textbf{0.918} & \textbf{0.658} & 6.57 & \textbf{0.683} & \textbf{0.744} \\
        \bottomrule
    \end{tabular}
    \end{scriptsize}
    }
\end{table*}


% \subsection{Spectral Transform-based Fine-tuning for Latent- and Pixel-Space Image Generation}
% \label{subsec:image_pixel}

\paragraph{Additional pixel-space fine-tuning results.} For pixel-space image generation, we evaluate our method on PixelGen-XXL/16 T2I~\citep{pixelgen} at its native resolution of $512^2$ pixels. We use the same fine-tuning data construction and LoRA training procedure as the Z-Image latent-space experiments above, with synthetic fine-tuning images generated by full-resolution PixelGen with 25 denoising steps from a set of 5K MS-COCO prompts disjoint from the 5K MS-COCO validation prompts used for evaluation. Since PixelGen is pretrained with $x_0$ prediction rather than flow matching, we convert the spectral-transform velocity targets from the fine-tuning procedure in Sec.~\ref{sec:method:lora} to $x_0$ targets and supervise the LoRA adapter with the converted $x_0$.

We use $S=2$ with $\delta=0.01$ and, following Sec.~\ref{subsec:image_latent}, organize the comparison into two wall-clock speedup tiers, all calculated relative to the 25-step PixelGen baseline on a single H100 GPU:

\begin{itemize}
    \item \textbf{Small speedup}: the 25-step native full-resolution PixelGen baseline, used as the reference for the reduced 13-step PixelGen variant and our training-free acceleration and LoRA fine-tuned variants at $S=2$ (25 steps), yielding approximately $1.5\times$ to $1.9\times$ speedup;
    \item \textbf{Medium speedup}: the 8-step native full-resolution PixelGen variant, used as an iso-compute reference for our training-free acceleration and LoRA fine-tuned variants at $S=2$ with reduced denoising steps, yielding approximately $2\times$ to $3\times$ speedup.
\end{itemize}

For quantitative evaluation, we follow the same protocol used in latent-space image generation and report the numbers in Table~\ref{tab:pixelgen-main}. We show extensive qualitative comparisons in Figures~\ref{fig:pixelgen-1}--\ref{fig:pixelgen-3}. Because PixelGen does not natively support a high-fidelity $256^2$ resolution generation, our training-free acceleration variant falls short of the reduced-step PixelGen baselines on overall image quality across both speedup tiers. Our spectral-transformation-based LoRA fine-tuning closes this gap: at the same FLOPs as the training-free acceleration variant, the LoRA fine-tuned model matches or surpasses the reduced-step PixelGen baseline on most image-quality and prompt-alignment metrics, including ImageReward, CLIP-IQA, and T2I-CompBench. The qualitative comparisons in Figures~\ref{fig:pixelgen-1}--\ref{fig:pixelgen-3} corroborate this trend, showing that our LoRA fine-tuned model's outputs match the visual quality of the 25-step native-resolution PixelGen baseline at substantially higher throughput.


\begin{table*}[t!]
    \centering
    \caption{\textbf{Quantitative comparisons on pixel-space image generation.} The number in parentheses next to PixelGen indicates the total number of inference steps. $\uparrow / \downarrow$ denotes that a higher / lower metric is favourable. Speedup is calculated relative to the 25-step PixelGen baseline wall-clock runtime on a single H100 GPU. \textcolor{blue}{\textbf{T}} and \textcolor{red}{\textbf{S}} denote temporal and spatial acceleration, respectively.}
    \label{tab:pixelgen-main}
    \resizebox{\textwidth}{!}{
    \begin{scriptsize}
    \begin{tabular}{cccccccccc}
         \toprule
        \multirow{2}{*}{Method} & \multirow{2}{*}{Accel.} & \multirow{2}{*}{Latency (s) $\downarrow$} & \multirow{2}{*}{Speedup (s) $\uparrow$} & \multirow{2}{*}{TFLOPs $\downarrow$} & Overall & \multicolumn{2}{c}{Image quality} & \multicolumn{2}{c}{Text alignment} \\
        \cmidrule(lr){6-10}
        & & & & & ImageReward $\uparrow$ & CLIP-IQA $\uparrow$ & NIQE $\downarrow$ & T2I-Comp. $\uparrow$ & GenEval $\uparrow$  \\
        \midrule
        PixelGen (25 steps) & - & 0.48 & 1.00$\times$ & 65.36 & 0.921 & 0.734 & 5.95 & 0.574 & 0.794 \\
        \cmidrule(lr){1-10}
        PixelGen (13 steps) & \textcolor{blue}{\textbf{T}} & \textbf{0.25} & \textbf{1.92$\times$} & \underline{34.16} & \underline{0.886} & \underline{0.726} & \underline{5.89} & \underline{0.571} & \underline{0.781} \\
        \textbf{Ours (TF, $\delta=0.01$, $S = 2$)}     & \textcolor{red}{\textbf{S}}       & \underline{0.30} & \underline{1.60$\times$} & \textbf{33.72} & 0.799 & 0.718 & 6.10 & 0.568 & \textbf{0.782} \\
        \textbf{Ours (LoRA, $\delta=0.01$, $S = 2$)}   & \textcolor{red}{\textbf{S}}       & 0.31 & 1.55$\times$ & \textbf{33.72} & \textbf{0.913} & \textbf{0.728} & \textbf{5.87} & \textbf{0.580} & 0.776 \\
        \cmidrule(lr){1-10}
        PixelGen (8 steps) & \textcolor{blue}{\textbf{T}} & \textbf{0.16} & \textbf{3.00$\times$} & \underline{21.16} & \underline{0.858} & \underline{0.715} & \textbf{5.79} & \underline{0.572} & 0.756 \\
        \textbf{Ours (TF, $\delta=0.01$, $S = 2$)}     & \textcolor{red}{\textbf{S}}       & \underline{0.21} & \underline{2.29$\times$} & \textbf{20.75} & 0.779 & 0.713 & 6.02 & 0.567 & \underline{0.765} \\
        \textbf{Ours (LoRA, $\delta=0.01$, $S = 2$)}   & \textcolor{red}{\textbf{S}}       & 0.22 & 2.18$\times$ & \textbf{20.75} & \textbf{0.908} & \textbf{0.724} & \underline{5.85} & \textbf{0.577} & \textbf{0.770} \\
        \bottomrule
    \end{tabular}
    \end{scriptsize}
    }
\end{table*}


\clearpage
\begin{figure}
    \centering
    \includegraphics[width=\linewidth]{supp_figures/flux_supp_1.pdf}
    \caption{\textbf{Qualitative comparisons on latent-space image generation.} We compare our method against default-step generation, reduced-step native-resolution generation on FLUX.1-dev \citep{black2024flux}, and RALU~\citep{ralu}, a state-of-the-art acceleration baseline matched to similar speedups. Our method outperforms both baselines. FLUX.1-dev with reduced steps degrades image quality and exhibits over-saturation artifacts, while RALU introduces noticeable noise artifacts.}
    \label{fig:flux-1}
\end{figure}

\clearpage
\begin{figure}
    \centering
    \includegraphics[width=\linewidth]{supp_figures/flux_supp_2.pdf}
    \caption{\textbf{Qualitative comparisons on latent-space image generation.} We compare our method against default-step generation, reduced-step native-resolution generation on FLUX.1-dev \citep{black2024flux}, and RALU~\citep{ralu}, a state-of-the-art acceleration baseline matched to similar speedups. Our method outperforms both baselines. FLUX.1-dev with reduced steps degrades image quality and exhibits over-saturation artifacts, while RALU introduces noticeable noise artifacts.}
    \label{fig:flux-2}
\end{figure}

\clearpage
\begin{figure}
    \centering
    \includegraphics[width=\linewidth]{supp_figures/flux_supp_3.pdf}
    \caption{\textbf{Qualitative comparisons on latent-space image generation.} We compare our method against default-step generation, reduced-step native-resolution generation on FLUX.1-dev \citep{black2024flux}, and RALU~\citep{ralu}, a state-of-the-art acceleration baseline matched to similar speedups. Our method outperforms both baselines. FLUX.1-dev with reduced steps degrades image quality and exhibits over-saturation artifacts, while RALU introduces noticeable noise artifacts.}
    \label{fig:flux-3}
\end{figure}

\iffalse
\clearpage
\begin{figure}
    \centering
    \includegraphics[width=\linewidth]{supp_figures/flux_supp_4.pdf}
    \caption{\textbf{Qualitative comparisons on latent-space image generation.} We compare our method against default-step generation, reduced-step native-resolution generation on FLUX.1-dev \citep{black2024flux}, and RALU~\citep{ralu}, a state-of-the-art acceleration baseline matched to similar speedups. Our method outperforms both baselines. FLUX.1-dev with reduced steps degrades image quality and exhibits over-saturation artifacts, while RALU introduces noticeable noise artifacts.}
    \label{fig:flux-4}
\end{figure}

\clearpage
\begin{figure}
    \centering
    \includegraphics[width=\linewidth]{supp_figures/flux_supp_5.pdf}
    \caption{\textbf{Qualitative comparisons on latent-space image generation.} We compare our method against default-step generation, reduced-step native-resolution generation on FLUX.1-dev \citep{black2024flux}, and RALU~\citep{ralu}, a state-of-the-art acceleration baseline matched to similar speedups. Our method outperforms both baselines. FLUX.1-dev with reduced steps degrades image quality and exhibits over-saturation artifacts, while RALU introduces noticeable noise artifacts.}
    \label{fig:flux-5}
\end{figure}
\fi

\clearpage
\begin{figure}
     \centering
     \includegraphics[width=\linewidth]{supp_figures/zimage_supp_1.pdf}
     \caption{\textbf{Qualitative comparisons on latent-space image generation (fine-tuned).} We compare our method against default-step generation, reduced-step native-resolution generation on Z-Image \citep{zimage} matched to similar speedups. Our fine-tuned model (Ours$^*$) achieves even higher image quality compared to our training-free acceleration variant and outperforms the reduced-step baseline.}
     \label{fig:zimage-1}
\end{figure}

\clearpage
\begin{figure}
     \centering
     \includegraphics[width=\linewidth]{supp_figures/zimage_supp_2.pdf}
     \caption{\textbf{Qualitative comparisons on latent-space image generation (fine-tuned).} We compare our method against default-step generation, reduced-step native-resolution generation on Z-Image \citep{zimage} matched to similar speedups. Our fine-tuned model (Ours$^*$) achieves even higher image quality compared to our training-free acceleration variant and outperforms the reduced-step baseline.}
     \label{fig:zimage-2}
\end{figure}

\clearpage
\begin{figure}
     \centering
     \includegraphics[width=\linewidth]{supp_figures/zimage_supp_3.pdf}
     \caption{\textbf{Qualitative comparisons on latent-space image generation (fine-tuned).} We compare our method against default-step generation, reduced-step native-resolution generation on Z-Image \citep{zimage} matched to similar speedups. Our fine-tuned model (Ours$^*$) achieves even higher image quality compared to our training-free acceleration variant and outperforms the reduced-step baseline.}
     \label{fig:zimage-3}
\end{figure}


\iffalse
\clearpage
\begin{figure}
     \centering
     \includegraphics[width=\linewidth]{supp_figures/zimage_supp_4.pdf}
     \caption{\textbf{Qualitative comparisons on latent-space image generation (fine-tuned).} We compare our method against default-step generation, reduced-step native-resolution generation on Z-Image \citep{zimage} matched to similar speedups. Our fine-tuned model (Ours$^*$) achieves even higher image quality compared to our training-free acceleration variant and outperforms the reduced-step baseline.}
     \label{fig:zimage-4}
\end{figure}

\clearpage
\begin{figure}
     \centering
     \includegraphics[width=\linewidth]{supp_figures/zimage_supp_5.pdf}
     \caption{\textbf{Qualitative comparisons on latent-space image generation (fine-tuned).} We compare our method against default-step generation, reduced-step native-resolution generation on Z-Image \citep{zimage} matched to similar speedups. Our fine-tuned model (Ours$^*$) achieves even higher image quality compared to our training-free acceleration variant and outperforms the reduced-step baseline.}
     \label{fig:zimage-5}
\end{figure}
\fi

\clearpage
\begin{figure}
    \centering
    \includegraphics[width=0.67\linewidth]{supp_figures/pixelgen_supp_1.pdf}
    \caption{\textbf{Qualitative comparisons on pixel-space image generation.} We compare our method against default-step generation and reduced-step native-resolution generation on PixelGen~\citep{pixelgen}, matched to comparable speedups. An asterisk (Ours$^*$) marks the fine-tuned model. Our method achieves similar quality to full-resolution generation while attaining higher speedups.}
    \label{fig:pixelgen-1}
\end{figure}

\clearpage
\begin{figure}
    \centering
    \includegraphics[width=0.67\linewidth]{supp_figures/pixelgen_supp_2.pdf}
    \caption{\textbf{Qualitative comparisons on pixel-space image generation.} We compare our method against default-step generation and reduced-step native-resolution generation on PixelGen~\citep{pixelgen}, matched to comparable speedups. An asterisk (Ours$^*$) marks the fine-tuned model. Our method achieves similar quality to full-resolution generation while attaining higher speedups.}
    \label{fig:pixelgen-2}
\end{figure}

\clearpage
\begin{figure}
    \centering
    \includegraphics[width=0.67\linewidth]{supp_figures/pixelgen_supp_3.pdf}
    \caption{\textbf{Qualitative comparisons on pixel-space image generation.} We compare our method against default-step generation and reduced-step native-resolution generation on PixelGen~\citep{pixelgen}, matched to comparable speedups. An asterisk (Ours$^*$) marks the fine-tuned model. Our method achieves similar quality to full-resolution generation while attaining higher speedups.}
    \label{fig:pixelgen-3}
\end{figure}

\iffalse
\clearpage
\begin{figure}
    \centering
    \includegraphics[width=0.67\linewidth]{supp_figures/pixelgen_supp_4.pdf}
    \caption{\textbf{Qualitative comparisons on pixel-space image generation.} We compare our method against default-step generation and reduced-step native-resolution generation on PixelGen~\citep{pixelgen}, matched to comparable speedups. An asterisk (Ours$^*$) marks the fine-tuned model. Our method achieves similar quality to reduced-step generation while attaining higher speedups.}
    \label{fig:pixelgen-4}
\end{figure}


\clearpage
\begin{figure}
    \centering
    \includegraphics[width=0.67\linewidth]{supp_figures/pixelgen_supp_5.pdf}
    \caption{\textbf{Qualitative comparisons on pixel-space image generation.} We compare our method against default-step generation and reduced-step native-resolution generation on PixelGen~\citep{pixelgen}, matched to comparable speedups. An asterisk (Ours$^*$) marks the fine-tuned model. Our method achieves similar quality to reduced-step generation while attaining higher speedups.}
    \label{fig:pixelgen-5}
\end{figure}
\fi


\clearpage
\section{Additional Video Generation Results}
\label{sec:video_generation}

\paragraph{Training-free acceleration of latent-space video generation.} For training-free inference acceleration of video generation, we use our internal cluster of NVIDIA H100 GPUs with 80~GB of VRAM. For quantitative evaluation, we generate 200 videos using 200 test prompts and report the standard VBench scores across six evaluation dimensions, including Subject Consistency, Background Consistency, Motion Smoothness, Dynamic Degree, Aesthetic Quality, and Image Quality, in Table~\ref{tab:wan_vbench_720p}. All videos are generated in 720P resolution. For additional video results and qualitative comparisons, please see the project website at~\url{https://howardxiao.ca/speed/}.


\clearpage
\section{Extended Ablation Studies}
\label{sec:ablation}

We perform extended ablation studies on latent-space image generation. Specifically, we vary the error tolerance $\delta$, the number of resolution stages $S$, and the spectral transformation $T_\Phi$, and report wall-clock latency, speedup, and ImageReward in Table~\ref{tab:ablation-quan}, with qualitative comparisons in Figures~\ref{fig:ablation-T} to \ref{fig:ablation-s}.

\begin{table*}[t!]
    \centering
    \caption{\textbf{Ablation studies on the spectral transformation $T_\Phi$, error tolerance $\delta$, and number of resolution stages $S$.} Speedup is calculated relative to the 7-step FLUX.1-dev baseline wall-clock runtime on a single H100 GPU.}
    \label{tab:ablation-quan}
    \centering
    \begin{tiny}
    \begin{tabular}{cccccc}
        \toprule
        $T_\Phi$ & $\delta$ & $S$ & Latency (s) $\downarrow$ & Speedup (s) $\uparrow$ & ImageReward $\uparrow$ \\
        \midrule
        FLUX (7 steps) & - & - & 2.08 & 1.00$\times$ & 0.920 \\
        \midrule
        DCT & 0.001 & 2 & \textbf{2.08} & \textbf{1.00$\times$} & 1.032 \\
        DWT & 0.001 & 2 & \underline{2.08} & \textbf{1.00$\times$} & \textbf{1.038} \\
        FFT & 0.001 & 2 & 2.08 & \textbf{1.00$\times$} & \underline{1.036} \\
        \midrule
        DCT & 0.001 & 3 & \textbf{2.05} & \textbf{1.01$\times$} & \underline{1.034} \\
        DWT & 0.001 & 3 & \underline{2.06} & \textbf{1.01$\times$} & \underline{1.034} \\
        FFT & 0.001 & 3 & 2.06 & \textbf{1.01$\times$} & \textbf{1.040} \\
        \midrule
        DCT & 0.001 & 4 & \textbf{2.05} & \textbf{1.02$\times$} & \textbf{1.008} \\
        DWT & 0.001 & 4 & \underline{2.05} & \underline{1.01$\times$} & \underline{1.002} \\
        FFT & 0.001 & 4 & 2.06 & \underline{1.01$\times$} & 0.992 \\
        \midrule
        DCT & 0.002 & 2 & \textbf{2.07} & \textbf{1.01$\times$} & \underline{1.032} \\
        DWT & 0.002 & 2 & \underline{2.08} & \underline{1.00$\times$} & 1.031 \\
        FFT & 0.002 & 2 & 2.08 & \underline{1.00$\times$} & \textbf{1.036} \\
        \midrule
        DCT & 0.002 & 3 & \textbf{2.05} & \textbf{1.02$\times$} & \underline{1.034} \\
        DWT & 0.002 & 3 & \underline{2.06} & \underline{1.01$\times$} & \underline{1.034} \\
        FFT & 0.002 & 3 & 2.06 & \underline{1.01$\times$} & \textbf{1.040} \\
        \midrule
        DCT & 0.002 & 4 & \textbf{2.05} & \textbf{1.02$\times$} & \textbf{1.008} \\
        DWT & 0.002 & 4 & \underline{2.05} & \underline{1.01$\times$} & \underline{1.002} \\
        FFT & 0.002 & 4 & 2.06 & \underline{1.01$\times$} & 0.992 \\
        \midrule
        DCT & 0.005 & 2 & \textbf{1.91} & \textbf{1.09$\times$} & \underline{1.019} \\
        DWT & 0.005 & 2 & \underline{1.92} & \textbf{1.09$\times$} & \textbf{1.026} \\
        FFT & 0.005 & 2 & 1.92 & \underline{1.08$\times$} & 1.014 \\
        \midrule
        DCT & 0.005 & 3 & \textbf{1.87} & \textbf{1.12$\times$} & \underline{1.009} \\
        DWT & 0.005 & 3 & \underline{1.88} & \underline{1.11$\times$} & \underline{1.009} \\
        FFT & 0.005 & 3 & 1.88 & \underline{1.11$\times$} & \textbf{1.028} \\
        \midrule
        DCT & 0.005 & 4 & \textbf{1.87} & \textbf{1.12$\times$} & \underline{1.005} \\
        DWT & 0.005 & 4 & \underline{1.87} & \underline{1.11$\times$} & \textbf{1.008} \\
        FFT & 0.005 & 4 & 1.87 & \underline{1.11$\times$} & 0.994 \\
        \midrule
        DCT & 0.01 & 2 & \textbf{1.75} & \textbf{1.19$\times$} & \textbf{1.039} \\
        DWT & 0.01 & 2 & \underline{1.76} & \underline{1.18$\times$} & \underline{1.025} \\
        FFT & 0.01 & 2 & 1.76 & \underline{1.18$\times$} & 1.016 \\
        \midrule
        DCT & 0.01 & 3 & \textbf{1.71} & \textbf{1.22$\times$} & 1.015 \\
        DWT & 0.01 & 3 & \underline{1.72} & \underline{1.21$\times$} & \underline{1.018} \\
        FFT & 0.01 & 3 & 1.72 & \underline{1.21$\times$} & \textbf{1.030} \\
        \midrule
        DCT & 0.01 & 4 & \textbf{1.71} & \textbf{1.22$\times$} & \textbf{1.014} \\
        DWT & 0.01 & 4 & \underline{1.71} & \textbf{1.22$\times$} & \underline{1.006} \\
        FFT & 0.01 & 4 & 1.72 & \underline{1.21$\times$} & 0.996 \\
        \midrule
        DCT & 0.02 & 2 & \textbf{1.60} & \textbf{1.30$\times$} & \underline{1.009} \\
        DWT & 0.02 & 2 & \underline{1.61} & \textbf{1.30$\times$} & \textbf{1.010} \\
        FFT & 0.02 & 2 & \underline{1.61} & \textbf{1.30$\times$} & 1.002 \\
        \midrule
        DCT & 0.02 & 3 & \textbf{1.53} & \textbf{1.36$\times$} & \underline{0.967} \\
        DWT & 0.02 & 3 & \underline{1.53} & \textbf{1.36$\times$} & 0.966 \\
        FFT & 0.02 & 3 & 1.53 & \textbf{1.36$\times$} & \textbf{0.980} \\
        \midrule
        DCT & 0.02 & 4 & \textbf{1.53} & \textbf{1.36$\times$} & \textbf{0.964} \\
        DWT & 0.02 & 4 & \underline{1.53} & \textbf{1.36$\times$} & \underline{0.955} \\
        FFT & 0.02 & 4 & 1.53 & \textbf{1.36$\times$} & 0.949 \\
        \midrule
        DCT & 0.05 & 2 & \textbf{1.44} & \textbf{1.44$\times$} & \underline{0.989} \\
        DWT & 0.05 & 2 & \underline{1.45} & \textbf{1.44$\times$} & \textbf{0.991} \\
        FFT & 0.05 & 2 & 1.45 & \underline{1.43$\times$} & 0.972 \\
        \midrule
        DCT & 0.05 & 3 & \textbf{1.34} & \textbf{1.55$\times$} & \underline{0.936} \\
        DWT & 0.05 & 3 & \underline{1.34} & \textbf{1.55$\times$} & \textbf{0.937} \\
        FFT & 0.05 & 3 & 1.35 & \underline{1.54$\times$} & 0.933 \\
        \midrule
        DCT & 0.05 & 4 & \textbf{1.33} & \textbf{1.56$\times$} & \textbf{0.918} \\
        DWT & 0.05 & 4 & \underline{1.34} & \textbf{1.56$\times$} & \underline{0.902} \\
        FFT & 0.05 & 4 & 1.34 & \underline{1.55$\times$} & 0.895 \\
        \midrule
        DCT & 0.1 & 2 & \textbf{1.29} & \textbf{1.62$\times$} & \underline{0.966} \\
        DWT & 0.1 & 2 & \underline{1.29} & \textbf{1.62$\times$} & \textbf{0.977} \\
        FFT & 0.1 & 2 & 1.30 & \underline{1.61$\times$} & 0.944 \\
        \midrule
        DCT & 0.1 & 3 & \textbf{1.15} & \textbf{1.81$\times$} & \underline{0.880} \\
        DWT & 0.1 & 3 & 1.17 & 1.78$\times$ & \textbf{0.892} \\
        FFT & 0.1 & 3 & \underline{1.17} & \underline{1.79$\times$} & 0.877 \\
        \midrule
        DCT & 0.1 & 4 & \textbf{1.15} & \textbf{1.82$\times$} & \underline{0.864} \\
        DWT & 0.1 & 4 & \underline{1.15} & \underline{1.81$\times$} & \textbf{0.867} \\
        FFT & 0.1 & 4 & 1.16 & 1.80$\times$ & 0.841 \\
        \bottomrule
    \end{tabular}
    \end{tiny}
\end{table*}


\paragraph{Computational efficiency and scaling.} The inference cost of Spectral Progressive Diffusion is determined by the number of denoising steps and the token count processed at each resolution stage. For a fixed pretrained model and resolution schedule, the per-sample cost is independent of the dataset size, so generating or evaluating more samples scales linearly with the number of samples. The spectral transformation and noise-expansion operations add only lightweight per-transition overhead compared to the DiT forward passes. 

\paragraph{Spectral Transformation $T_\Phi$.} In Fig. \ref{fig:ablation-T}, we observe that DCT, DWT, and FFT achieve similar high-fidelity performance. We default to DCT for its real-valued computation and native support of non-power-of-two resolution scaling.

\paragraph{Error tolerance $\delta$.} In Fig. \ref{fig:ablation-delta}, we observe that increasing $\delta$ improves efficiency, but results in ghosting and halo artifacts near detailed edges.

\paragraph{Resolution stages $S$.} In Fig. \ref{fig:ablation-s}, we find that increasing $S$ leads to marginal speedup improvements and little image quality degradation. Since the primary speedup derives from bypassing native high-resolution steps, $S=2$ achieves the largest speedup while higher $S$ provides diminishing efficiency gains by further optimizing the already-inexpensive low-resolution trajectory.




\newpage
\begin{figure}
    \centering
    \includegraphics[width=\linewidth]{supp_figures/ablation_T.pdf}
    \caption{\textbf{Qualitative ablation on $T_\Phi$.} We see that all spectral transformations attain similar image quality. }
    \label{fig:ablation-T}
\end{figure}

\newpage
\begin{figure}
    \centering
    \includegraphics[width=\linewidth]{supp_figures/ablation_delta.pdf}
    \caption{\textbf{Qualitative ablation on $\delta$.} We observe that increasing $\delta$ improves efficiency, but results in ghosting and halo artifacts near detailed edges. }
    \label{fig:ablation-delta}
\end{figure}

\newpage
\begin{figure}
    \centering
    \includegraphics[width=\linewidth]{supp_figures/ablation_s.pdf}
    \caption{\textbf{Qualitative ablation on $S$.} We find that increasing $S$ leads to marginal speedup improvements and little image quality degradation.}
    \label{fig:ablation-s}
\end{figure}


\clearpage
\section{Additional Frequency-based Image Editing Details}
\label{sec:image_editing}

\paragraph{Detailed Procedure of Frequency-based Image Editing.} In our frequency-based editing pipeline, we perform the following operations given an input image $\mathbf{x}_\text{in}$, resolution scales $s_{1:S}$, and transition times $t_{1:S-1}$. We assume that the image editing process starts from transition time $t_k$, with $k \in \{1, 2, \ldots, S-1\}$.

\begin{enumerate}
    \item Compute the spectrum of the input image $\boldsymbol{\xi}_\text{in} = T_\Phi(\mathbf{x}_\text{in})$.
    \item Extract the low-frequency part of $\boldsymbol{\xi}_\text{in}$ corresponding to the representable spectrum $\Omega_{s_k}$ at scale $s_k$. 
    \item Fill the spectrum $\Omega_{s_{k+1}} \setminus \Omega_{s_k}$ with $t_k\epsilon$, where $\epsilon \sim \mathcal{N}(0, 1)$.
    \item Convert the expanded spectrum $\Omega_{s_{k+1}}$ back to the spatial domain $\mathbf{x}_{t_k}^{s_{k+1}}$ via $T^{-1}_\Phi$ and continue denoising starting from the transition time $t_k$.
\end{enumerate}

The total time of denoising is thus $1-t_k$.

We conduct image editing experiments using Z-Image \citep{zimage} as the backbone, evaluating our frequency-domain image editing pipeline against the SDEdit baseline \citep{meng2022sdedit}. We set $S=2$ and  $\delta=0.01$.

\paragraph{Extended Qualitative Results.} We show extended image editing results on two tasks: \textit{texture editing} and \textit{artistic stylization}. We configure our method with $\delta=0.01$ and $S=2$, which corresponds to a resolution transition at step 26 of the default 50-step denoising trajectory. We set $T_\Phi=\text{DCT}$, and evaluate the effect of different spectral transforms in the next section. For the SDEdit baseline, we initialize the process by corrupting the clean source image to the equivalent noise level at step 26 and subsequently proceed with the denoising schedule. This setting accurately reflects the editing method's capability to inject textural or style information, which are considered high-frequency details, in later denoising steps. 

As shown in Figures \ref{fig:texture_edit_1} to \ref{fig:art_stylization}, our frequency-based editing method achieves superior prompt alignment and geometric consistency compared to SDEdit, successfully changing the texture and artistic style of the input image while retaining the topological structure of the input. While SDEdit-style editing at earlier timesteps (higher noise levels at step 17) can facilitate stylistic changes, it introduces significant structural drift and geometric divergence from the source image. Furthermore, starting from early timesteps incurs substantial computational overhead and requires a brittle renoising to a specific, difficult-to-tune initialization timestep to balance editability with structural fidelity.

\paragraph{Effect of $T_\Phi$ on Image Editing.} As shown in Fig. \ref{fig:T_comparison_editing}, FFT-, DCT- and DWT-based editing achieve similar editing quality.

\begin{figure}
    \centering
    \includegraphics[width=\linewidth]{supp_figures/image_editing_supp_1.pdf}
    \caption{\textbf{Texture editing results.} Our frequency-based editing framework outperforms SDEdit, enabling high-fidelity texture transfer while preserving the geometric structure of the input image.}
    \label{fig:texture_edit_1}
\end{figure}

\begin{figure}
    \centering
    \includegraphics[width=\linewidth]{supp_figures/image_editing_supp_2.pdf}
    \caption{\textbf{Texture editing results.} Our frequency-based editing framework outperforms SDEdit, enabling high-fidelity texture transfer while preserving the geometric structure of the input image.}
    \label{fig:texture_edit_2}
\end{figure}

\begin{figure}
    \centering
    \includegraphics[width=\linewidth]{supp_figures/image_editing_supp_3.pdf}
    \caption{\textbf{Texture editing results.} Our frequency-based editing framework outperforms SDEdit, enabling high-fidelity texture transfer while preserving the geometric structure of the input image.}
    \label{fig:texture_edit_3}
\end{figure}

\begin{figure}
    \centering
    \includegraphics[width=\linewidth]{supp_figures/image_editing_supp_4.pdf}
    \caption{\textbf{Texture editing results.} Our frequency-based editing framework outperforms SDEdit, enabling high-fidelity texture transfer while preserving the geometric structure of the input image.}
    \label{fig:texture_edit_4}
\end{figure}


\begin{figure}
    \centering
    \includegraphics[width=\linewidth]{supp_figures/artistic_supp.pdf}
    \caption{\textbf{Artistic stylization results.} Aside from texture editing, our frequency-based editing approach also supports artistic stylization given stylistic descriptions and a representative artist. }
    \label{fig:art_stylization}
\end{figure}

\begin{figure}
    \centering
    \includegraphics[width=\linewidth]{supp_figures/image_editing_T_comparison.pdf}
    \caption{\textbf{Effect of $T_\Phi$ on image editing.} FFT-, DCT- and DWT-based editing achieve similar editing quality.}
    \label{fig:T_comparison_editing}
\end{figure}


\clearpage
\section{Broader Impacts}
\label{app:broader-impacts}

Spectral Progressive Diffusion is a principled method for improving the efficiency of pretrained image and video generation models, rather than a system targeted at a particular deployment domain. Its main positive impact is to reduce the compute, latency, and energy cost of generation, which can make experimentation and creative workflows more accessible and lower the environmental footprint of repeated generation. At the same time, because the method can accelerate existing generative models, we acknowledge that it could also lower the cost of producing synthetic media for undesirable uses such as disinformation, impersonation, or spam, depending on the pretrained model and release setting. We do not release new datasets or generation services, and any future release of fine-tuned adapters should follow the access terms, safeguards, and usage restrictions of the underlying pretrained models.

\section{Existing Assets and Licenses}
\label{app:assets-licenses}

Our experiments build on publicly available pretrained models, datasets, evaluation benchmarks, and software packages, including FLUX.1-dev~\citep{black2024flux}, Z-Image~\citep{zimage}, PixelGen~\citep{pixelgen}, WAN~\citep{wan}, MS-COCO~\citep{lin2014microsoft}, GenEval~\citep{geneval}, T2I-CompBench~\citep{t2icompbench}, VBench~\citep{vbench}, and the evaluation metrics used in Sec.~\ref{sec:exp}, including ImageReward~\citep{imagereward}, CLIP-IQA~\citep{wang2023clipiqa}, and NIQE~\citep{mittal2013niqe}. We cite the original sources for these assets and use them under their respective licenses and access terms.
The named licenses and access terms are: FLUX.1-dev is released under the FLUX.1 [dev] Non-Commercial License; Z-Image, PixelGen, WAN~2.1, VBench, ImageReward, VChitect-T2V-Dataverse, and the PyIQA implementation used for NIQE are under the Apache License 2.0; GenEval and T2I-CompBench are under the MIT License; and CLIP-IQA is under the NTU S-Lab License 1.0.
MS-COCO annotations are under Creative Commons Attribution 4.0 International (CC BY 4.0), while its images follow their original Flickr licenses; the public page for Aesthetic-Train-V2 does not list an explicit license, so we use it only under its stated access terms and do not redistribute the dataset.

% \subsection{Training-free Acceleration of Latent-space Video Generation}
% \label{app:video_generation}

% We run all latent-space video generation experiments on a single H100 GPU and measure wall-clock time measurements. For quantitative evaluation, we generate 200 videos using 200 test prompts and report the standard VBench scores across six evaluation dimensions, including Subject Consistency, Background Consistency, Motion Smoothness, Dynamic Degree, Aesthetic Quality, and Image Quality, in Table~\ref{tab:wan_vbench_720p}. All videos are generated in 720P resolution. For additional video results and qualitative comparisons, please see the supplementary website (\texttt{website/index.html}).

% \section{Classification of general progressive-resolution transforms}
% \label{app:family}
% \hx{Under construction...}

% Progressive resolution generation requires a principled transition between a low-resolution flow state and a high-resolution flow state. In this appendix we characterize \emph{all} such transitions in closed form: they form an affine-linear family determined by an upsample operator $U$ and an associated noise-injection matrix $B$. Three specific choices of $U$ (block-repeat, bilinear, sinc) recover Haar wavelet, spatial interpolation, and Fourier-domain constructions respectively.

% \subsection{Problem Setup}
% \label{app:family:setup}

% Let $x \in \mathbb{R}^{N_l}$ be a clean data sample at low resolution, and let $y \in \mathbb{R}^{N_h}$ denote the corresponding clean sample at high resolution, related by a fixed linear upsample $U \in \mathbb{R}^{N_h \times N_l}$ with $y = U x$ and $N_h > N_l$. Write $k = N_h / N_l$ for the total count ratio (so $k = c^d$ if each of $d$ spatial dimensions grows by a factor $c$).

% At some diffusion time $t \in (0, 1)$, the low-resolution rectified-flow state is
% \begin{equation}
%   x_t \;=\; (1-t)\, x + t\, \epsilon_l,
%   \qquad \epsilon_l \sim \mathcal{N}(0, I_{N_l}),
%   \label{eq:low-res-state}
% \end{equation}
% and we seek a \emph{resolution lift} $L_t : \mathbb{R}^{N_l} \to \mathbb{R}^{N_h}$ (possibly using auxiliary randomness $\xi$) whose output is a valid flow state of $y = Ux$ at some target time $t'$:
% \begin{equation}
%   L_t(x_t;\, \xi)
%   \;\stackrel{d}{=}\;
%   (1-t')\, U x + t'\, \epsilon_h,
%   \qquad \epsilon_h \sim \mathcal{N}(0, I_{N_h}),
%   \label{eq:high-res-state}
% \end{equation}
% with $\xi$ independent of the low-resolution noise $\epsilon_l$. The requirement that $\epsilon_h$ be i.i.d.\ standard Gaussian is essential: the pretrained denoiser $v_\theta$ was trained on precisely this distribution, and any departure produces out-of-distribution inputs.

% \subsection{Characterization: $L_t$ Must Be Affine-Linear}
% \label{app:family:affine}

% \begin{theorem}[Affine-linearity]
% \label{thm:affine}
% Any $L_t$ satisfying \eqref{eq:high-res-state} for all clean signals $x \in \mathbb{R}^{N_l}$ and all realizations of $\epsilon_l$ is an affine-linear function of its first argument:
% \begin{equation}
%   L_t(y; \xi) \;=\; A\, y + B\, \xi,
%   \label{eq:affine}
% \end{equation}
% where $A \in \mathbb{R}^{N_h \times N_l}$, $B \in \mathbb{R}^{N_h \times m}$ for some $m$, and $\xi \sim \mathcal{N}(0, I_m)$ is an auxiliary Gaussian independent of the argument.
% \end{theorem}

% \begin{proof}[Proof sketch]
% The right-hand side of \eqref{eq:high-res-state} is Gaussian with mean $(1-t')\, U x$ \emph{linear} in $x$ and covariance $t'^2\, I_{N_h}$ \emph{independent of $x$}. Translating $x \mapsto x + \Delta x$ shifts the argument of $L_t$ by $(1-t)\, \Delta x$ and must shift the output distribution by $(1-t')\, U \Delta x$. Since $\xi$ is independent of the argument, this shift property holds pointwise in $\xi$:
% \begin{equation*}
%   L_t\bigl(y + (1-t)\,\Delta x;\; \xi\bigr)
%   \;=\;
%   L_t(y;\; \xi) \,+\, (1-t')\, U\, \Delta x .
% \end{equation*}
% A measurable function satisfying Cauchy's functional equation $f(y + \Delta) - f(y) = \ell(\Delta)$ for all $\Delta$ must be affine-linear in $y$; this yields $L_t(y; \xi) = A\, y + g(\xi)$ for some linear $A$ and measurable $g$. Finally, $L_t(y; \xi)$ must be Gaussian for every fixed $y$ (so that the distributional identity in~\eqref{eq:high-res-state} is preserved under conditioning on $y$), which forces $g(\xi) = B\, \xi$ for a linear $B$.
% \end{proof}

% \subsection{Closed-Form Solution for $A$ and $B$}
% \label{app:family:closedform}

% Given Theorem~\ref{thm:affine}, we substitute $y = x_t = (1-t)\, x + t\, \epsilon_l$ into \eqref{eq:affine} and match moments against the target distribution in \eqref{eq:high-res-state}.

% \paragraph{Mean matching.} Linearity in $x$ gives $A\, (1-t) x = (1-t')\, U x$ for all $x$, so
% \begin{equation}
%   \boxed{\;\;A \;=\; \frac{1-t'}{1-t}\, U\;\;}
%   \label{eq:A-closedform}
% \end{equation}
% The signal-path operator is therefore the upsample $U$, rescaled by the ratio of signal coefficients at the two times.

% \paragraph{Covariance matching.} The output covariance is the sum of contributions from $\epsilon_l$ (via $A$) and from $\xi$ (via $B$), and must equal $t'^2\, I_{N_h}$:
% \begin{equation}
%   t^2\, A A^\top + B B^\top \;=\; t'^2\, I_{N_h}.
%   \label{eq:cov-match}
% \end{equation}
% Substituting \eqref{eq:A-closedform} and solving for $B B^\top$,
% \begin{equation}
%   \boxed{\;\;B B^\top \;=\; t'^2\, I_{N_h} \,-\, t^2 \!\left(\tfrac{1-t'}{1-t}\right)^{\!2}\, U U^\top\;\;}
%   \label{eq:B-closedform}
% \end{equation}
% Any $B$ whose Gram matrix equals the right-hand side of \eqref{eq:B-closedform} (e.g., its symmetric square root) realizes a valid lift.

% \paragraph{$B$ is necessarily nonzero.} Since $U \in \mathbb{R}^{N_h \times N_l}$ maps between spaces of different dimension, $\operatorname{rank}(U U^\top) \le N_l < N_h$. The target $t'^2 I_{N_h}$ has full rank $N_h$. Hence the right-hand side of \eqref{eq:B-closedform} has rank at least $N_h - N_l > 0$: \emph{any genuine resolution lift requires injecting fresh Gaussian noise}, whose variance fills the $N_h - N_l$ directions not reached by $U$. This is a direct structural consequence of the dimension increase and is independent of the specific choice of $U$.

% \paragraph{PSD feasibility and the minimum-$t'$ theorem.} For $B$ to exist as a real matrix, $B B^\top$ must be positive semi-definite. Applying this to \eqref{eq:B-closedform},
% \begin{equation}
%   t'^2\, I_{N_h} \;\succeq\; t^2\! \left(\tfrac{1-t'}{1-t}\right)^{\!2} U U^\top,
%   \label{eq:PSD-condition}
% \end{equation}
% which reduces to a scalar inequality on the maximum eigenvalue of $U U^\top$:
% \begin{equation}
%   t'^2 \;\ge\; t^2\!\left(\tfrac{1-t'}{1-t}\right)^{\!2} \lambda_{\max}(U U^\top).
%   \label{eq:PSD-scalar}
% \end{equation}
% This bound is \emph{always active} for DC-preserving upsamples, as we now show.

% \begin{theorem}[Minimum-$t'$ bound for DC-preserving upsamples]
% \label{thm:min-t-prime}
% Let $U \in \mathbb{R}^{N_h \times N_l}$ be a DC-preserving upsample, i.e., $U\,\mathbf{1}_{N_l} = \mathbf{1}_{N_h}$ where $\mathbf{1}_n \in \mathbb{R}^n$ is the all-ones vector. Let $k = N_h / N_l$. Then $\lambda_{\max}(U U^\top) \ge k$, and the PSD feasibility condition \eqref{eq:PSD-scalar} combined with $t, t' \in (0, 1)$ forces
% \begin{equation}
%   \boxed{\;\; t' \;\ge\; \frac{\sqrt{\lambda_{\max}(U U^\top)} \; t}{1 + \bigl(\sqrt{\lambda_{\max}(U U^\top)} - 1\bigr)\, t} \;\;}
%   \label{eq:min-t-prime}
% \end{equation}
% with equality achievable when $\lambda_{\max}(U U^\top) = k$ (i.e., when $U$ attains the DC-preserving lower bound). In particular, for any $U$ that saturates $\lambda_{\max}(U U^\top) = k$,
% \begin{equation}
%   t'_{\min} \;=\; \frac{\sqrt{k}\, t}{1 + (\sqrt{k} - 1)\, t},
%   \label{eq:min-t-prime-tight}
% \end{equation}
% giving the tightest possible target time while preserving flow-matching structure.
% \end{theorem}

% \begin{proof}[Proof sketch]
% The Rayleigh quotient at the all-ones vector gives $\lambda_{\max}(U U^\top) \ge \|U\,\mathbf{1}_{N_l}\|_2^2 / \|\mathbf{1}_{N_l}\|_2^2 = \|\mathbf{1}_{N_h}\|_2^2 / N_l = N_h/N_l = k$. Let $s := \sqrt{\lambda_{\max}(U U^\top)}$. Taking square roots of \eqref{eq:PSD-scalar} (all quantities positive for $t, t' \in (0,1)$), $t'(1-t) \ge s\, t (1-t')$. Rearranging, $t' + s\,t\,t' \ge s\,t + t\,t'$, so $t'\bigl(1 + (s-1)\,t\bigr) \ge s\,t$, yielding \eqref{eq:min-t-prime}. Equality in \eqref{eq:PSD-scalar} is achieved when the eigenspace of $\lambda_{\max}$ saturates and all other eigenvalues lie below $\lambda_{\max}$; this is exactly the condition that $U U^\top$ has a rank-one eigenvector at the all-ones direction with eigenvalue $k$ and no direction exceeds $k$, i.e., $\|U\|_{\mathrm{op}}^2 = k$.
% \end{proof}

% \emph{Interpretation.} Equation \eqref{eq:min-t-prime-tight} captures a fundamental ``rewind'' required at every resolution transition: since $t'_{\min} > t$ for any $k > 1$, the target state must always sit at a \emph{noisier} time than the source. The rewind grows with $k$: at $t = 0.5$, $t'_{\min} = \sqrt{k}/(\sqrt{k}+1)$ yields $0.67$ for $k{=}4$, $0.75$ for $k{=}9$, and $0.80$ for $k{=}16$. Upsamples that amplify a non-DC direction above $k$ (e.g., bicubic, which exhibits a spectral overshoot) carry $\lambda_{\max} > k$ and therefore pay a strictly larger rewind than block-repeat, bilinear, and sinc---all of which saturate $\lambda_{\max}(U U^\top) = k$ exactly and thus achieve the tight minimum \eqref{eq:min-t-prime-tight}.

% \subsection{Three Canonical Instances of $U$}
% \label{app:family:instances}

% We now specialize \eqref{eq:A-closedform}--\eqref{eq:B-closedform} to three choices of $U$ used throughout the literature. We give closed-form expressions for $A$ and $B B^\top$ in each case, and illustrate them in the elementary $2 \!\times\! 2 \to 4 \!\times\! 4$ example (so $N_l = 4$, $N_h = 16$, and $k = 4$).

% \paragraph{(i) Block-repeat upsample (Haar LL-only IDWT).}
% The nearest-neighbor lift replicates each low-resolution coefficient $x_{(i)}$ to the $k$ high-resolution positions it covers:
% \begin{equation}
%   [U_{\mathrm{BR}}\, x]_{(i,\alpha)} \;=\; x_{(i)},
%   \qquad \alpha \in \{0, \ldots, k-1\}.
%   \label{eq:BR-def}
% \end{equation}
% For the $2\!\times\!2 \to 4\!\times\!4$ case, each input pixel is replicated into its $2\!\times\!2$ output block. The Gram matrix $U_{\mathrm{BR}} U_{\mathrm{BR}}^\top$ is block-diagonal: each $N_l$-group of $k$ replicated positions yields a $k \!\times\! k$ all-ones block $J_k$, with zero off-block entries. Its eigendecomposition has eigenvalue $k$ on the per-block constant vector and $0$ on the $(k{-}1)$-dimensional orthogonal complement (per block). Substituting into \eqref{eq:A-closedform}--\eqref{eq:B-closedform},
% \begin{equation}
%   A_{\mathrm{BR}} = \tfrac{1-t'}{1-t}\, U_{\mathrm{BR}},
%   \qquad
%   B_{\mathrm{BR}} B_{\mathrm{BR}}^\top = t'^2\, I_{N_h} \;-\; t^2\!\left(\tfrac{1-t'}{1-t}\right)^{\!2}\,
%   \bigoplus_{i=1}^{N_l} J_k .
%   \label{eq:BR-AB}
% \end{equation}
% Equivalently, $B_{\mathrm{BR}} B_{\mathrm{BR}}^\top$ is block-diagonal with per-block matrix $t'^2 I_k - t^2\!\left(\tfrac{1-t'}{1-t}\right)^{\!2} J_k$, whose eigenvalues are $t'^2$ on the $(k{-}1)$-dim complement of the constant direction and $t'^2 - k\, t^2\!\left(\tfrac{1-t'}{1-t}\right)^{\!2}$ on the constant direction itself.

% \emph{Connection to the Haar wavelet.} $U_{\mathrm{BR}}$ is exactly $\sqrt{k}$ times the orthonormal Haar inverse DWT applied to the coefficient vector $(x, 0, \ldots, 0)$ (low-frequency slot populated; detail slots zero). The block-diagonal decomposition into \emph{per-block constant direction} and \emph{per-block detail complement} is precisely the LL-vs-\{LH, HL, HH\} split of the Haar basis in the $2\!\times\!2 \to 4\!\times\!4$ case.

% \paragraph{(ii) Bilinear upsample.}
% The bilinear lift convolves $x$ with a separable tent kernel, yielding weighted averages of neighboring coefficients. For the $2\!\times\!2 \to 4\!\times\!4$ case with half-pixel offset, each output pixel receives four weights $w_1 w_1', w_1 w_2', w_2 w_1', w_2 w_2'$ with $(w_1, w_2) = (\tfrac{3}{4}, \tfrac{1}{4})$ in each dimension:
% \begin{equation}
%   [U_{\mathrm{BL}}\, x]_{(2i+a, 2j+b)}
%   \;=\; \sum_{p,q \in \{-1,0\}} w_{a-p}\, w_{b-q}\, x_{(i+p, j+q)}.
%   \label{eq:BL-def}
% \end{equation}
% The Gram matrix $U_{\mathrm{BL}} U_{\mathrm{BL}}^\top$ is now a \emph{banded, non-block-diagonal} matrix: adjacent $2\!\times\!2$ blocks share weights, so correlation leaks across block boundaries. It has the same maximum eigenvalue $k$ (attained on the all-ones direction, since $U_{\mathrm{BL}}$ preserves DC by construction $\sum_\alpha w_\alpha = 1$), but its eigenstructure decays smoothly with frequency rather than collapsing to $\{k, 0\}$. Substituting:
% \begin{equation}
%   A_{\mathrm{BL}} = \tfrac{1-t'}{1-t}\, U_{\mathrm{BL}},
%   \qquad
%   B_{\mathrm{BL}} B_{\mathrm{BL}}^\top = t'^2\, I_{N_h} \,-\, t^2\!\left(\tfrac{1-t'}{1-t}\right)^{\!2}\, U_{\mathrm{BL}} U_{\mathrm{BL}}^\top.
%   \label{eq:BL-AB}
% \end{equation}
% Taking the symmetric square root gives a frequency-shaped noise injector that is strongest where $U_{\mathrm{BL}} U_{\mathrm{BL}}^\top$ has the smallest eigenvalues (high spatial frequencies, where the bilinear transfer function rolls off). Unlike the block-repeat case, $B_{\mathrm{BL}}$ has no exact spatial locality.

% \paragraph{(iii) Sinc (ideal bandlimited) upsample.}
% In the Fourier domain, the ideal lift zero-pads the spectrum: $U_{\mathrm{S}}$ takes $x$, computes its orthonormal DFT $X = \mathcal{F}_{N_l} x$, places $X$ in the low-frequency block of an $N_h$-point spectrum, and inverts. Algebraically,
% \begin{equation}
%   U_{\mathrm{S}} \;=\; \sqrt{k}\, \mathcal{F}_{N_h}^{-1}\, P_L\, \widetilde{\mathcal{F}}_{N_l},
%   \qquad
%   P_L : \mathbb{R}^{N_h} \to \mathbb{R}^{N_h}
%   \text{ is the projection onto the low-frequency support},
%   \label{eq:sinc-def}
% \end{equation}
% where $\widetilde{\mathcal{F}}_{N_l}$ embeds the $N_l$-point spectrum into the low-frequency block of an $N_h$-point grid. For $2\!\times\!2 \to 4\!\times\!4$, $P_L$ selects the two lowest horizontal and two lowest vertical frequency bins (the central $2\!\times\!2$ block of the $4\!\times\!4$ shifted spectrum). The Gram matrix is a \emph{pure projection}:
% \begin{equation}
%   U_{\mathrm{S}} U_{\mathrm{S}}^\top \;=\; k\, P_L,
%   \label{eq:sinc-gram}
% \end{equation}
% with eigenvalues $k$ on the $N_l$-dim low-frequency subspace and $0$ on the $N_h - N_l$-dim high-frequency complement. Substituting,
% \begin{equation}
%   A_{\mathrm{S}} = \tfrac{1-t'}{1-t}\, U_{\mathrm{S}},
%   \qquad
%   B_{\mathrm{S}} B_{\mathrm{S}}^\top = t'^2\, I_{N_h} \,-\, k\, t^2\!\left(\tfrac{1-t'}{1-t}\right)^{\!2}\, P_L
%   \;=\; \underbrace{\Big(t'^2 - k\, t^2\!\left(\tfrac{1-t'}{1-t}\right)^{\!2}\Big)\, P_L}_{\text{low-freq residual}} \;+\; t'^2\, P_H,
%   \label{eq:sinc-AB}
% \end{equation}
% where $P_H = I_{N_h} - P_L$. When the rewind time $t'$ is chosen so that the low-frequency coefficient vanishes (i.e., $t' = \sqrt{k}\, t / (1 + (\sqrt{k}-1) t)$), $B_{\mathrm{S}}$ reduces to a \emph{pure high-pass} noise injector: $B_{\mathrm{S}} = t'\, P_H$, which simply fills the high-frequency half of the $N_h$-spectrum with fresh i.i.d.\ complex Gaussian noise. This is the cleanest form of \eqref{eq:B-closedform} and motivates the Fourier interpretation of the lift.

% \emph{Connection to FFT padding.} The sinc lift is exactly ``FFT, zero-pad spectrum, inverse FFT'' when the detail-band noise is zero. With the prescribed $B_{\mathrm{S}}$, it becomes ``FFT, pad with complex Gaussian noise at amplitude $t'$, inverse FFT''---a single-operator drop-in replacement for the wavelet lift.

% \paragraph{Summary.}
% All three instances share the same $A = \tfrac{1-t'}{1-t}\, U$ in \eqref{eq:A-closedform} and saturate the DC eigenvalue $\lambda_{\max}(U U^\top) = k$; they differ only in how $B B^\top$ distributes the required $N_h - N_l$ units of fresh variance. The block-repeat choice concentrates $B$ in per-block detail directions (Haar-local), bilinear spreads $B$ as a smooth high-pass filter (mildly non-local), and sinc projects $B$ onto the exact high-frequency subspace of the global DFT (fully non-local, frequency-ideal).



% \iffalse
% \section{Additional Quantitative Results}
% \label{app:additional-quant}

% \subsection{SD3 Main Results}
% \label{app:sd3-results}

% Table~\ref{tab:sd3-supp} replicates the main-paper FLUX.1-dev table (Table~\ref{tab:flux-main}) on Stable Diffusion~3~\citep{sd3} using the same RALU~\citep{ralu} evaluation protocol, at SD3's native 28 steps and 1024$\times$1024 resolution. Speedup targets are 2$\times$ and 3$\times$ per RALU. RAS~\citep{ras} is included as an SD3-specific baseline from the RALU table.

% \begin{table}[t]
%   \centering
%   \footnotesize
%   \caption{\textbf{Image results on Stable Diffusion~3}, matching RALU's~\citep{ralu} SD3 table. Placeholder values.}
%   \label{tab:sd3-supp}
%   \setlength{\tabcolsep}{2.5pt}
%   \begin{tabular}{l ccc ccccc}
%     \toprule
%     Method & Latency $\downarrow$ & TFLOPs $\downarrow$ & Speedup $\uparrow$ & FID $\downarrow$ & NIQE $\downarrow$ & CLIP-IQA $\uparrow$ & T2I-Comp.\ $\uparrow$ & GenEval $\uparrow$ \\
%     \midrule
%     SD3 (28 steps, base)                   & - & - & 1$\times$ & - & - & - & - & - \\
%     SD3 (reduced steps)                    & - & - & -         & - & - & - & - & - \\
%     ToCa~\citep{toca}                      & - & - & -         & - & - & - & - & - \\
%     Bottleneck Sampling~\citep{bottleneck} & - & - & -         & - & - & - & - & - \\
%     RAS~\citep{ras}                        & - & - & -         & - & - & - & - & - \\
%     RALU~\citep{ralu}                      & - & - & -         & - & - & - & - & - \\
%     Ours (spectral passthrough)            & - & - & -         & - & - & - & - & - \\
%     \textbf{Ours-TF (Haar)}                & - & - & -         & - & - & - & - & - \\
%     \textbf{Ours-TF (FFT)}                 & - & - & -         & - & - & - & - & - \\
%     \textbf{Ours-LoRA (Haar)}              & - & - & -         & - & - & - & - & - \\
%     \bottomrule
%   \end{tabular}
% \end{table}

% \subsection{PixelGen (Pixel-Space) Main Results}
% \label{app:pixelgen-results}

% Table~\ref{tab:pixelgen-supp} reports the same protocol on the pixel-space backbone PixelGen~\citep{pixelgen}, which operates directly on RGB pixels rather than a VAE latent.

% \begin{table}[t]
%   \centering
%   \footnotesize
%   \caption{\textbf{Pixel-space image results on PixelGen}~\citep{pixelgen}. Same RALU-style columns as Table~\ref{tab:flux-main}; baselines restricted to methods applicable in pixel space. Placeholder values.}
%   \label{tab:pixelgen-supp}
%   \setlength{\tabcolsep}{2.5pt}
%   \begin{tabular}{l ccc ccccc}
%     \toprule
%     Method & Latency $\downarrow$ & TFLOPs $\downarrow$ & Speedup $\uparrow$ & FID $\downarrow$ & NIQE $\downarrow$ & CLIP-IQA $\uparrow$ & T2I-Comp.\ $\uparrow$ & GenEval $\uparrow$ \\
%     \midrule
%     PixelGen (base)                        & - & - & 1$\times$ & - & - & - & - & - \\
%     PixelGen (reduced steps)               & - & - & -         & - & - & - & - & - \\
%     Bottleneck Sampling~\citep{bottleneck} & - & - & -         & - & - & - & - & - \\
%     Ours (spectral passthrough)            & - & - & -         & - & - & - & - & - \\
%     \textbf{Ours-TF (Haar)}                & - & - & -         & - & - & - & - & - \\
%     \textbf{Ours-TF (FFT)}                 & - & - & -         & - & - & - & - & - \\
%     \textbf{Ours-LoRA (Haar)}              & - & - & -         & - & - & - & - & - \\
%     \bottomrule
%   \end{tabular}
% \end{table}

% \fi
